\pdfoutput=1
\documentclass{amsart}

%==============================================================================
% PACKAGES
%==============================================================================

%------------------------------------------------------------------------------

%----- encoding, fonts & language ---------------------------------------------
\usepackage[T1]{fontenc}
%\usepackage[utf8]{inputenc}
\usepackage[mathscr]{eucal}% so-called Euler fonts, allows \mathscr
\usepackage{bbold}% allows \unit \mathbb{1}
\usepackage{microtype}

%----- math packages ----------------------------------------------------------
\usepackage{amsmath}% allows pmatrix
\usepackage{amssymb}% allows special arrows like >-> e.g.
\usepackage{amsthm}% allows theorem* , e.g.
\usepackage{comment}
\usepackage{mathtools}
\usepackage{centernot}% allows \centernot\Longrightarrow for better "not" placement
\usepackage[a]{esvect}

%----- colours & graphics -----------------------------------------------------
\usepackage[usenames,dvipsnames]{xcolor} %changed to xcolor; color was giving ``incompatible color definition'' warnings
\usepackage{graphicx}
\usepackage{adjustbox}
\usepackage{pict2e}
\usepackage[most]{tcolorbox}
%\usepackage{color}

%----- diagrams ---------------------------------------------------------------
\usepackage{wasysym}% allows LEFTCIRCLE
\usepackage{tikz}
%\usepackage{tkz-berge}
\usepackage{tikz-cd}
\usetikzlibrary{arrows}
\usetikzlibrary{decorations.markings}

\usepackage[all]{xy}
%\xyoption{color}
%\xyoption{pdf}
\xyoption{line}
%\CompileMatrices
\usepackage{quiver}

%----- tables, lists & text formatting ----------------------------------------
\usepackage{booktabs} %pretty tables
\usepackage{makecell}
\usepackage{longtable}
\usepackage{tablefootnote}
\usepackage{arydshln}
\usepackage[normalem]{ulem}% allows \sout to strike-out, cross-out text.
%\usepackage{enumerate}% allows easy change of label
\usepackage{enumitem}% allows resuming enumerate
%\usepackage{array}% allows array
\usepackage{caption}
\captionsetup{belowskip=-10pt}%reduces the vspace after a figure
\usepackage{pdflscape}

%----- miscallaneous & utility ------------------------------------------------
\usepackage{comment}% allows multiple lines to be commeted out
\usepackage[perpage]{footmisc} %resets footnotes per page
\usepackage{etoolbox} %this is needed for the hack removing address indentation

%----- collaboration & annotations --------------------------------------------
%\usepackage[disable]{todonotes}
\usepackage{todonotes}
\reversemarginpar % puts the todonotes on the right so they don't clash with showkeys; doesn't seem to be working

%some other colours: Violet, Red, Brown, ...
\newcommand{\Beren}[1]{{\color{MidnightBlue}#1}}
\newcommand{\Drew}[1]{{\color{Orange}#1}}
\newcommand{\Tobi}[1]{{\color{Green}#1}}
\newcommand{\Natalia}[1]{{\color{Yellow}#1}}
\newcommand{\Dout}[1]{\Drew{\sout{#1}}}
\newcommand{\Bout}[1]{\Beren{\sout{#1}}}
\newcommand{\Tout}[1]{\Tobi{\sout{#1}}}
\newcommand{\Nout}[1]{\Natalia{\sout{#1}}}
\newcommand{\Greg}[1]{{\color{magenta}#1}}

\newcommand{\ttodo}[1]{\todo[color=green!40]{#1}}
\newcommand{\tcomment}[1]{\todo[inline,color=green!40]{#1}}
\newcommand{\btodo}[1]{\todo[color=blue!20]{#1}}
\newcommand{\bcomment}[1]{\todo[inline,color=blue!20]{#1}}
\newcommand{\ntodo}[1]{\todo[color=yellow!40]{#1}}
\newcommand{\ncomment}[1]{\todo[inline,color=yellow!40]{#1}}
\newcommand{\dtodo}[1]{\todo[color=Orange!40]{#1}}
\newcommand{\dcomment}[1]{\todo[inline,color=Orange!40]{#1}}
\newcommand{\gtodo}[1]{\todo[color=magenta!40]{#1}}
\newcommand{\gcomment}[1]{\todo[inline,color=magenta!40]{#1}}
\newcommand{\ctodo}[1]{\todo[color=pink!40]{#1}}
\newcommand{\ccomment}[1]{\todo[inline,color=pink!40]{#1}}

%==============================================================================
% HYPERLINKS, LABELS & CROSS-REFERENCING
%==============================================================================
% note: hyperref should generally be loaded late, followed by cleveref.
%==============================================================================

% get rid of following for final version
%\usepackage[left]{showlabels} %showlabels isn't good with left equation labels
%\usepackage[notref,notcite]{showkeys}% fait apparaitre les references
%\usepackage[hyphenbreaks]{breakurl} %not needed when using pdflatex

%----- hyperref --------------------------------------------

%\usepackage{hyperref}
%To generate  D V I  with links using latex
%\usepackage[hypertex]{hyperref} % SEE BELOW: Comment this one out and uncomment the one below when you create PDF.
%to generate  P D F  with links using latex
\usepackage[unicode]{hyperref} %added unicode; was getting warnings otherwise

%\hypersetup{colorlinks=true,citecolor=Brown,urlcolor=Brown,linkcolor=Brown}% quote away to generate ps% replace BrickRed by Brown?

\definecolor{dark-red}{rgb}{0.5,0.15,0.15}
\definecolor{dark-blue}{rgb}{0.15,0.15,0.6}
\definecolor{dark-green}{rgb}{0.15,0.6,0.15}

%\hypersetup{
%    colorlinks, linkcolor=OliveGreen,
%    citecolor=OliveGreen, urlcolor=OliveGreen
%}
\hypersetup{
    colorlinks, linkcolor=Blue,
    citecolor=Blue, urlcolor=Blue
}

%----- cleveref --------------------------------------------

\usepackage[nameinlink,capitalise,noabbrev]{cleveref}

%==============================================================================
% SETTINGS & VARIABLES
%==============================================================================

% to mark overful hbox's; alternative to using [draft] option because that deactivates the hyperlinks
\overfullrule=5mm

%\hyphenation{right-equivariant}

\numberwithin{equation}{section}% makes equat numb contain the section
\setcounter{tocdepth}{1}% to have only sections in the T.o.C.
\renewcommand{\labelenumi}{(\alph{enumi})}% makes enumeration via a, b, c...
\renewcommand{\labelenumii}{(\alph{enumi}.\arabic{enumii})}% idem for subcases


%%todo: is this for booktabs?? when is this used?
\setlength\dashlinedash{0.2pt}
\setlength\dashlinegap{1.5pt}
\setlength\arrayrulewidth{0.3pt}

%==============================================================================
% LOW-LEVEL HACKS, TIKZ & XY-PIC CONFIGURATIONS
%==============================================================================

%----- XY-pic configurations --------------------------------------------------
\newdir{ >}{{}*!/-10pt/\dir{>}}

%----- TikZ styles ------------------------------------------------------------
\tikzset{
    labelrotatebelow/.style={anchor=north, rotate=90, inner sep=1.0mm}
}
\tikzset{
    labelrotateabove/.style={anchor=south, rotate=90, inner sep=1.0mm}
}

\tikzset{negated/.style={
        decoration={markings,
            mark= at position 0.5 with {
                \node[transform shape] (tempnode) {$\backslash \! \! \backslash$};
            }
        },
        postaction={decorate}
    }
}

%----- low-level LaTeX macros & hacks -----------------------------------------
\makeatletter
\providecommand*{\twoheadrightarrowfill@}{%
  \arrowfill@\relbar\relbar\twoheadrightarrow
}
\providecommand*{\twoheadleftarrowfill@}{%
  \arrowfill@\twoheadleftarrow\relbar\relbar
}
\providecommand*{\xtwoheadrightarrow}[2][]{%
  \ext@arrow 0579\twoheadrightarrowfill@{#1}{#2}%
}
\providecommand*{\xtwoheadleftarrow}[2][]{%
  \ext@arrow 5097\twoheadleftarrowfill@{#1}{#2}%
}
\makeatother

%%%%hardcore hacks
\makeatletter
\newcommand{\cveewedge@measure}[2]{%
  \sbox\z@{$#1\m@th#2$}%
  \dimen@=1.05\ht\z@
  \unitlength=.005\wd\z@
  \count@=\dimen@ 
  \divide\count@\unitlength
  \ifx#1\scriptstyle
    \linethickness{0.8\@wholewidth}%
  \else
    \ifx#1\scriptscriptstyle
      \linethickness{0.65\@wholewidth}%
    \fi
  \fi
}

\newcommand{\cwedge}{\mathbin{\mathpalette\do@cwedge\relax}}
\newcommand{\do@cwedge}[2]{%
  \cveewedge@measure{#1}{\wedge}
  \begin{picture}(200,\count@)
  \roundjoin
  \polygon(25,0)(100,\count@)(175,0)
  \end{picture}%
}
\newcommand{\cvee}{\mathbin{\mathpalette\do@cvee\relax}}
\newcommand{\do@cvee}[2]{%
  \cveewedge@measure{#1}{\vee}
  \begin{picture}(200,\count@)
  \roundjoin
  \polygon(25,\count@)(100,0)(175,\count@)
  \end{picture}%
}
\makeatother
%%%%more hardcore hacks
\makeatletter
\newcommand{\leqnomode}{\tagsleft@true\let\veqno\@@leqno}
\newcommand{\reqnomode}{\tagsleft@false\let\veqno\@@eqno}
\makeatother
%%%%end hardcore hacks

%==============================================================================
% THEOREM ENVIRONMENTS & CREF NAMES
%==============================================================================

%---------- other theorem-environments packages -----
%---------- label related packages -----
%---------- label related packages -----
%---------- math packages -----
\newtheorem{ThmAlpha}{Theorem}
\renewcommand{\theThmAlpha}{\Alph{ThmAlpha}}
\swapnumbers %pour que le numéro apparaisse devant le théorème

\newtheorem{Thm}[equation]{Theorem}
%\newtheorem{MainThm}[equation]{Main Theorem}
\newtheorem*{Thm*}{Theorem}
\newtheorem*{MainThm*}{Main Theorem}
%\newtheorem{ThmDef}[equation]{Theorem and Definition}
\newtheorem{Prop}[equation]{Proposition}
%\newtheorem{PropDef}[equation]{Proposition and Definition}
\newtheorem{Lem}[equation]{Lemma}
%\newtheorem{MainLem}[equation]{Main Lemma}
\newtheorem{Cor}[equation]{Corollary}
\newtheorem{CorConj}[equation]{Corollary of Conjecture}
%\newtheorem*{Cor*}{Corollary}
\newtheorem{Conj}[equation]{Conjecture}
\newtheorem{Leap}[equation]{Leap of Faith}
%\newtheorem{Axiom}[equation]{Axiom}
\newtheorem*{Que*}{Question}
\newtheorem*{Goal*}{Goal}
\newtheorem{Claim}[equation]{Claim}
\newtheorem{Hope}[equation]{Hope}

%\theoremstyle{definition}
\theoremstyle{remark}
\newtheorem{Def}[equation]{Definition}
\newtheorem{Ter}[equation]{Terminology}
%\newtheorem{Defs}[equation]{Definitions}
\newtheorem{Not}[equation]{Notation}
%\newtheorem{Nots}[equation]{Notations}
%\newtheorem{DefNot}[equation]{Definition-Notation}
\newtheorem{Exa}[equation]{Example}
\newtheorem*{Exau}{Example}
\newtheorem{Exas}[equation]{Examples}
\newtheorem{Cons}[equation]{Construction}
\newtheorem{Conv}[equation]{Convention}
\newtheorem{Hyp}[equation]{Hypothesis}

%\theoremstyle{remark}
\newtheorem{Rem}[equation]{Remark}
\newtheorem{Par}[equation]{}
\newtheorem{Rec}[equation]{Recollection}
%\newtheorem{Rems}[equation]{Remarks}
\newtheorem{Que}[equation]{Question}
%\newtheorem{Baksh}[equation]{Baksheesh}
\newtheorem{War}[equation]{Warning}

\Crefname{Thm}{Theorem}{Theorems}
\Crefname{Prop}{Proposition}{Propositions}
\Crefname{Lem}{Lemma}{Lemmas}
\Crefname{Cor}{Corollary}{Corollaries}
\Crefname{Exa}{Example}{Examples}
\Crefname{ThmAlpha}{Theorem}{Theorems}
\Crefname{Rem}{Remark}{Remarks}
\Crefname{Rec}{Recollection}{Recollections}
\Crefname{Def}{Definition}{Definitions}
\Crefname{War}{Warning}{Warnings}
\Crefname{Cons}{Construction}{Constructions}
\newcommand{\crefrangeconjunction}{--}

%==============================================================================
% COMMANDS & MACROS
%==============================================================================

\newcounter{enum-resume-hack}
\renewcommand{\thepart}{\Roman{part}}


%----- math macros & operators -----
\newcommand{\nc}{\newcommand}
\nc{\dmo}{\DeclareMathOperator}

\renewcommand{\emptyset}{\varnothing}
%\renewcommand{\complement}{{\scriptscriptstyle\mathsf{C}}}
\renewcommand{\complement}{{\mkern-1mu\mathsf{C}}}
\renewcommand{\complement}{{\mathsf{C}}}

\newcommand{\eqv}{\mathrel{\makebox[\widthof{$\Longleftrightarrow$}][c]{$\Longleftrightarrow$}}}
\newcommand{\imp}{\mathrel{\makebox[\widthof{$\Longleftrightarrow$}][c]{$\Longrightarrow$}}}

\nc{\Rpq}{R_{[\frakp,\frakq]}}
\dmo{\Soc}{Soc}
\dmo{\rank}{rank}
\dmo{\Iso}{Iso}
\dmo{\CBheight}{ht_{\hspace{-0.1ex}\text{\tiny CB}}\hspace{-0.10ex}}
\nc{\CBdegree}{\deg_{\hspace{-0.1ex}\text{\tiny CB}}\hspace{-0.10ex}}
\nc{\Rp}{R_{\frakp}}
\nc{\Rq}{R_{\frakq}}
\dmo{\kos}{kos}
\dmo{\BL}{BL}
%\nc{\EB}{E_{\beta}}
\nc{\GammaR}{\Gamma_{\hspace{-0.15ex}R}}
\nc{\KtimesRtimes}{K^\times\hspace{-0.3ex}/R^\times}
\nc{\ABbar}{\bar{\cat A}_{\cat B}}
\nc{\ACbar}{\bar{\cat A}_{\cat C}}
\nc{\EBbar}{\bar{E}_{\cat B}}
\nc{\ECbar}{\bar{E}_{\cat C}}
\nc{\unitbar}{\bar{\unit}}
\nc{\Fhat}{\hat{F}}
\nc{\Uhat}{\hat{U}}
\nc{\Fbar}{\bar{F}}
\nc{\Ubar}{\bar{U}}
\nc{\EB}{E_{\cat B}}
\nc{\EC}{E_{\cat C}}
\nc{\varphiiC}{\varphi_i^h\hspace{-0.1ex}(\cat C)}
\nc{\varphiC}{\varphi^h\hspace{-0.1ex}(\cat C)}
\nc{\varphiB}{\varphi^h\hspace{-0.1ex}(\cat B)}
\nc{\EvarphiiC}{E_{\varphiiC}}
\nc{\EvarphiC}{E_{\varphiC}}
\nc{\EvarphiB}{E_{\varphiB}}
\nc{\GammaA}{\Gamma_{\hspace{-0.4ex}A}}
\nc{\LA}{L_{\hspace{-0.1ex}A}}
\nc{\LB}{L_{\hspace{-0.1ex}B}}
\nc{\LeT}{L_e\cat T}
\nc{\LfT}{L_{\hspace{-0.1ex}f}\cat T}
\nc{\LAT}{\LA\cat T}
\nc{\LBT}{\LB\cat T}
\nc{\LAS}{\LA\cat S}
\nc{\LBS}{\LB\cat S}
\nc{\LKBT}{L_{\KB}\cat T}
\nc{\KB}{K(\cat B)}
\nc{\KS}{K\hspace{-0.2ex}(S)}
\nc{\abcX}{a\hspace{-0.3ex}\bc{X}}
\nc{\abcx}{a\hspace{-0.3ex}\bc{x}}
\nc{\meet}{\wedge}%\curlywedge
\nc{\overbar}[1]{\mkern 1.5mu\overline{\mkern-1.5mu#1\mkern-1.5mu}\mkern 1.5mu}
\renewcommand{\L}{\mathrm{L}}
%\nc{\BCdual}[1]{\Delta_{\bbQ/\bbZ}(#1)}
\nc{\weaklyfinite}{weakly closed}
\nc{\finite}{closed}
\nc{\BCdual}[1]{{#1}^*}
\nc{\LCore}{\mathrm{LCore}}
\nc{\Stovicek}{{\v{S}\v{t}ov\'{i}\v{c}ek}}
\nc{\ftriple}{f_{\natural}}
\nc{\unitC}{\unit_{\cat C}}%
\nc{\unitD}{\unit_{\cat D}}%
%\nc{\Pone}{\mathbb{P}^1_{\hspace{-0.2em}k}}
\nc{\Pone}{{\mathbb{P}^1}}
\nc{\InvSupp}[1]{\Supp_{\cat T}^{-1}(#1)}
\nc{\InvCosupp}[1]{\Cosupp_{\cat T}^{-1}(#1)}
\nc{\closureP}{\overbar{\{\cat P\}}}
\nc{\closureQ}{\overbar{\{\cat Q\}}}
\nc{\singP}{\{\cat P\}}
\nc{\singQ}{\{\cat Q\}}
\nc{\singm}{\{\frak m\}}
\dmo{\Inj}{Inj}
\dmo{\tfib}{tfib}
\dmo{\tcof}{tcof}
\dmo{\Aut}{Aut}
\dmo{\tofib}{tofib}
\dmo{\surj}{surj}
\dmo{\Excs}{Exc}
\dmo{\Homog}{Homog}
\dmo{\PSh}{PSh}
\dmo{\Epis}{Epi}
\newcommand{\Exc}[1]{\Excs_{#1}}
\newcommand{\Epi}[1]{\Epis_{\le #1}}
\dmo{\KInjdmo}{K}
\dmo{\Dbdmo}{mod}
\dmo{\sur}{sur}
\nc{\KInj}[1]{\KInjdmo(\Inj #1)}
\nc{\Dbmod}[1]{\Der^b(\Dbdmo #1)}
\dmo{\Viss}{vis}%lol
\nc{\Vis}{\Viss}
\nc{\vis}{\Vis}
\nc{\kappaaux}{g}
\nc{\kappaCh}{{\kappaaux(\cat C_h)}}
\nc{\kappam}{{\kappaaux({\frak m})}}
\nc{\kappaP}{{\kappaaux_{\cat P}}}
\nc{\kappaQ}{{\kappaaux(\cat Q)}}
\nc{\kappaCP}{{\kappaaux_{\cat C}(\cat P)}}
\nc{\kappaDP}{{\kappaaux_{\cat D}(\cat P)}}
\nc{\kappaCQ}{{\kappaaux_{\cat C}(\cat Q)}}
\nc{\kappaDQ}{{\kappaaux_{\cat D}(\cat Q)}}
\nc{\kappaphiB}{{\kappaaux(\phi(\cat B))}}
\nc{\kappaphiQ}{{\kappaaux(\varphi(\cat Q))}}
\nc{\halfplus}{{\scriptscriptstyle\top}}
\dmo{\Sub}{Sub}
%\nc{\SpEn}{\Sp_{L_n}}
%\nc{\SpEn}{\SH_{L_n}}
\nc{\SpEn}{\cat S_{E(n)}}
\nc{\SpEnf}{\cat S_n}
%\nc{\SpEnf}{\SH_{L_n^f}}
%\nc{\SpEnf}{\Sp_{L_n^f}}
\nc{\Lcomp}{L^{\mathrm{com}}} %I made this and the next one commands because I'm unsure of the choice of notation
\nc{\Ucomp}{U^{\mathrm{com}}}
\nc{\bbullet}{{\scriptscriptstyle\hspace{-1pt}\bullet}}
\nc{\bullett}{{\scriptscriptstyle\bullet}\hspace{-1pt}}
\nc{\LF}{L\hspace{-0.2ex}F}
\dmo{\StMod}{StMod}
\dmo{\Proj}{Proj}
\nc{\SpG}{\Sp^G}
\nc{\EG}{\bbE_G}
\nc{\DEG}{\Der(\EG)}
\nc{\DE}{\Der(\bbE)}
\nc{\Prst}{{\cat P}\mathrm{r^{st}}}
%\nc{\Mack}{\mathcal{M}ack}
\nc{\Mack}{\mathrm{Mack}}
%\nc{\DerKal}{\Der \hspace{-0.2em} \mathcal{M}_{\mathrm{Kal}}}
\nc{\SC}{S\cat C}
\dmo{\fin}{{fin}}
\dmo{\DM}{DM}
\dmo{\fp}{fp}
\nc{\DMQ}{\DM_Q}
%\dmo{\SpecMack}{SpectralMack}
\dmo{\DerKal}{DMack}
\dmo{\coh}{coh}
\dmo{\Der}{D}
\dmo{\DMot}{DMot}
\dmo{\Cell}{Cell}
\dmo{\rmH}{H}
\dmo{\piu}{\underline{\pi}}
\dmo{\Sphere}{\mathbb{S}}
\nc{\HA}{{\rmH \hspace{-0.2em}\bbA}}
\nc{\HZ}{{\rmH \hspace{-0.2em}\bbZ}}
\nc{\HZbar}{{\rmH \hspace{-0.2em}\underline{\bbZ}}}
\nc{\HbbF}{{\rmH \hspace{-0.15em}\mathbb{F}}}
\nc{\Fp}{{\bbF_{\hspace{-0.1em}p}}}
\nc{\HFp}{{\rmH \hspace{-0.15em}\bbF_{\hspace{-0.1em}p}}}
\nc{\HQ}{{\rm H \bbQ}}
\nc{\DHZG}{\Der(\HZ_G)}
\nc{\DHZH}{\Der(\HZ_H)}
\nc{\DHZK}{\Der(\HZ_K)}
\nc{\DHZGN}{\Der(\HZ_{G/N})}
\nc{\DHZGG}{\Der(\HZ_{G/G})}
\nc{\DHZCp}{\Der(\HZ_{C_p})}
\nc{\DHZGprime}{\Der(\HZ_{G'})}
\nc{\DHZ}{\Der(\HZ)}
\nc{\frakr}{\mathfrak{r}}
\nc{\frakp}{\mathfrak{p}}
\nc{\frakq}{\mathfrak{q}}
\nc{\frakS}{\mathfrak{S}}
\nc{\frakT}{\mathfrak{T}}
\nc{\Z}{\mathbb{Z}}
\nc{\F}{\mathbb{F}}
\nc{\SSG}{\text{sSet}_*^G}
\nc{\sSet}{\text{sSet}}
\newcommand{\cosuppa}{\operatorname{co-supp}}

\dmo{\csupp}{csupp}
\dmo{\Con}{Conj}
\dmo{\Id}{Id}
\dmo{\rmK}{\textrm{\rm K}}
\dmo{\Spc}{Spc}
\dmo{\thick}{thick}
\dmo{\thickid}{thickid}
\nc{\thicko}[1]{\thickid\langle #1 \rangle}
\nc{\thickt}[1]{\thick_\otimes\langle #1 \rangle}
\dmo{\cone}{cone}
\dmo{\End}{End}
\dmo{\Derperf}{D_{perf}}
\dmo{\Mor}{Mor}
\dmo{\id}{id}
\dmo{\incl}{incl}
\dmo{\Img}{Im}
\dmo{\im}{im}
\dmo{\Ker}{Ker}
\dmo{\ind}{ind}
\dmo{\CoInd}{coind}
\dmo{\GH}{GH}
\dmo{\idem}{e}
\dmo{\res}{res}
\dmo{\infl}{infl}
\dmo{\Derqc}{D_{qc}}
\nc{\DbcohX}{\Der^b(\coh X)}
\dmo{\triv}{triv}
\dmo{\Pic}{Pic}
\dmo{\dual}{dual}
\dmo{\Tel}{Tel} %telescope
\dmo{\grMod}{grMod}%
\dmo{\Mod}{Mod}%
%\dmo{\mod}{mod}%
\dmo{\opname}{op}
\dmo{\SH}{SH}% ground name for cat of spectra
\dmo{\smallb}{b}% ground exponent for ``bounded''
\dmo{\Spec}{Spec}
\dmo{\supp}{supp}
\dmo{\Supp}{Supp}
\dmo{\crosseffec}{cr}
\newcommand{\crosseffect}{{\crosseffec}}
\dmo{\thofib}{tofib}
\dmo{\hofib}{hofib}
\dmo{\cosupp}{cosupp}
\dmo{\Cosupp}{Cosupp}
\nc{\SHc}{{\SH^c}}
\nc{\SHp}{{\SH_{(p)}}}
\nc{\SHcp}{{\SH^c_{(p)}}}
\nc{\SHG}{\SH(G)}
\nc{\SHGp}{\SH(G)_{(p)}}
\nc{\SHGc}{\SHG^c}
\nc{\SHGcp}{\SHG^c_{(p)}}
\nc{\quadtext}[1]{\quad\textrm{#1}\quad}
\nc{\qquadtext}[1]{\qquad\textrm{#1}\qquad}
\nc{\adj}{\dashv}
\nc{\adjto}{\rightleftarrows}
\nc{\bbL}{\mathbb{L}}
\nc{\bbS}{\mathbb{S}}
\nc{\bbA}{\mathbb{A}}
\nc{\bbE}{\mathbb{E}}
\nc{\bbN}{\mathbb{N}}
\nc{\bbQ}{\mathbb{Q}}
\nc{\bbZ}{\mathbb{Z}}
\nc{\bbR}{\mathbb{R}}
\nc{\bbF}{\mathbb{F}}
\nc{\cat}[1]{\mathscr{#1}}%or: \nc{\cat}[1]{\mathcal{#1}}
\nc{\ie}{{\sl i.e.}, }
\nc{\into}{\mathop{\rightarrowtail}}
\nc{\inv}{^{-1}}
\nc{\isoto}{\mathop{\overset{\sim}\to}}
\nc{\isotoo}{\mathop{\overset{\sim}\too}}
\nc{\onto}{\mathop{\twoheadrightarrow}}
\nc{\too}{\mathop{\longrightarrow}\limits}
\nc{\mapstoo}{\longmapsto}
\nc{\adh}[1]{\overline{#1}}% adherence
\nc{\adhpt}[1]{\adh{\{#1\}}}% adherence of a pt
\nc{\aka}{{a.\,k.\,a.}\ }
\nc{\calF}{\mathcal{F}}
\nc{\eg}{{\sl e.\,g.}}
\nc{\hook}{\hookrightarrow}

\nc{\ideal}[1]{\langle #1\rangle}
\dmo{\red}{red}
\dmo{\Hom}{Hom}
\nc{\Homcat}[1]{\Hom_{\cat #1}}
%\nc{\ihom}{{\underline{\hom}}}
\nc{\iHom}{\mathcal{H}\mathrm{om}}
\nc{\Supph}{\Supp^h}
\nc{\Supphnaive}{\Supp^n}
\nc{\Cosupph}{\Cosupp^h}
\nc{\ihom}[1]{\mathsf{hom}(#1)}
\nc{\ihomC}[1]{\mathsf{hom}_{\cat C}(#1)}
\nc{\ihomD}[1]{\mathsf{hom}_{\cat D}(#1)}
\nc{\ihomsub}[2]{\mathsf{hom}_{#1}(#2)}
%\nc{\ihom}[1]{F(#1)}
%\nc{\ihom}[1]{\mathrm{hom}(#1)}
%\nc{\ihom}[1]{\mathcal{H}\mathrm{om}(#1)}
\nc{\Mid}{\,\big|\,}
\nc{\MMod}{\,\text{-}\Mod}%
\nc{\GrMMod}{\,\text{-}\grMod}%
\nc{\op}{^{\opname}}
\nc{\oto}[1]{\overset{#1}\to}
\nc{\otoo}[1]{\overset{#1}{\,\too\,}}
\nc{\sminus}{\!\smallsetminus\!}
\nc{\poplus}[1]{^{\oplus #1}}%
\nc{\potimes}[1]{^{\otimes #1}}% tensor power
\nc{\sbull}{{\scriptscriptstyle\bullet}}%\mathbf{\cdot}}%{}}
\nc{\SET}[2]{\left\{\,#1\middle\vert#2\,\right\}}
\nc{\SETT}[2]{\left\{\,#1\,\middle\vert\,#2\,\right\}}
\nc{\SpcK}{\Spc(\cat K)}% most used
\nc{\then}{\Rightarrow}
\nc{\unit}{\mathbb{1}}% unit for \otimes
\nc{\xra}{\xrightarrow}
\nc{\phigeom}[1]{\widetilde{\Phi}^{#1}}
\dmo{\Oname}{O}
\dmo{\proper}{proper}% for proper subgroups
\dmo{\lenormal}{\unlhd}
\dmo{\fib}{fib}
\dmo{\cofib}{cofib}
\dmo{\lnormal}{\lhd}
\nc{\normal}{\trianglelefteq}%\lhd
\nc{\Op}{\Oname^p}% O^p for maximal p-normal subgroup
\nc{\Oq}{\Oname^q}% as above for p=q
\newcommand{\Sp}{{\mathscr{S}p}}
\newcommand{\Sppc}{{\mathscr{S}p^{\raisebox{1pt}{$\scriptstyle c$}}_{\hspace{-0.2ex}(p)}}}
\newcommand{\Spp}{{\mathscr{S}p_{\hspace{-0.2ex}(p)}}}
\newcommand{\Top}{{\mathscr{T}op}}
\dmo{\Ho}{Ho}
\dmo{\CB}{CB}
\dmo{\Fin}{Fin}
\dmo{\add}{add}
%\dmo{\Add}{Add}
\dmo{\Fun}{Fun}
\dmo{\Ext}{Ext}
\dmo{\Tor}{Tor}
\dmo{\ann}{ann}
\dmo{\Ann}{Ann}
\dmo{\CAlg}{CAlg}
\dmo{\CMon}{CMon}
\dmo{\CC}{\cat C} 
\dmo{\DD}{\cat D}
\dmo{\OO}{\mathcal{O}}
\dmo{\Map}{Map}
\dmo{\Span}{Span}
\dmo{\Tot}{Tot}
\dmo{\N}{N}
\dmo{\Cat}{Cat}
\dmo{\colim}{colim}
\dmo{\hocolim}{hocolim}
\dmo{\Ch}{Ch}
\dmo{\A}{\mathbb{A}^{eff}}
\nc{\AGeff}{\mathbb{A}_G^{\mathrm{eff}}}
\nc{\BGeff}{\mathcal{B}_G^{\mathrm{eff}}}
\nc{\BG}{{\mathcal{B}_G}}
\nc{\NBGeff}{{\N}{\BGeff}}
%\nc{\Ab}{\mathbf{Ab}}
\dmo{\Ab}{Ab}
\nc{\Smith}{\mathsf{Smith}}
\nc{\Floyd}{\mathsf{Floyd}}
\nc{\blue}{\beth^{\mathrm{geom}}}
\dmo{\Set}{Set}
\dmo{\ev}{ev}
\dmo{\Spcl}{Spcl}
\nc{\Funadd}{\Fun_{\add}}
\dmo{\proj}{proj}
\dmo{\cof}{cof}
%\dmo{\ind}{ind}
\nc{\cPd}{\cat P_{\hspace{-.1em}d}}
\nc{\cPm}{\cat P_{\hspace{-.1em}m}}
\nc{\Chp}{\mathsf{Ch}_p}
\nc{\Pp}{\mathsf{P}_p}


%New commands
\dmo{\Coideal}{Coideal}
\dmo{\gen}{gen}
\nc{\auxcoidealsymb}{\vartriangleleft}

\dmo{\Loc}{Loc}
\dmo{\Ind}{Ind}
\dmo{\Coloc}{Coloc}
\dmo{\Locideal}{Locid}
\dmo{\Colocideal}{Coloc}
\nc{\LOCO}{\Locideal}
\nc{\COLOCO}{\Colocideal}
\dmo{\Kos}{Kos}
\nc{\Loco}[1]{\LOCO\langle #1 \rangle}
\nc{\Coloco}[1]{\COLOCO\langle #1 \rangle}

\nc{\LambdaP}{\Lambda^{\cat P}}
\nc{\LambdaQ}{\Lambda^{\cat Q}}
\nc{\GammaP}{\Gamma_{\cat P}} 
\nc{\GammaQ}{\Gamma_{\cat Q}} 
\nc{\LambdaW}{\Lambda^{\hspace{-0.3ex}W}}
\nc{\GammaW}{\Gamma_{\hspace{-0.3ex}W}}
\nc{\gW}{g_W}
\nc{\gP}{g_{\cat P}}
\nc{\gQ}{g_{\cat Q}}
\nc{\cC}{{\cat C}}
\nc{\cT}{{\cat T}}
\nc{\cD}{{\cat D}}

\nc{\mT}{\kern-0.5em\mod\kern-0.1em\text{-}\cat{T}^c}
\nc{\mTc}{\kern-0.5em\mod\kern-0.1em\text{-}\cat{T}^c}
\nc{\MTc}{\Mod\kern-0.1em\text{-}\cat{T}^c}
\nc{\MT}{\Mod\kern-0.1em\text{-}\cat{T}}

\DeclareMathOperator{\type}{type}
\DeclareMathOperator{\height}{height}
\DeclareMathOperator{\pos}{pos}


\DeclareMathOperator{\coind}{coind}

\DeclareMathOperator{\map}{map}

\nc{\cL}{\mathcal{L}}
\nc{\cP}{\mathcal{P}}

\nc{\tblue}{\beth^{\mathrm{Tate}}}


\nc{\cA}{\mathcal{A}}
\nc{\cF}{\mathcal{F}}
\nc{\Fnt}{\cF_{\mathrm{nt}}}
\nc{\Fall}{\cF_{\mathrm{all}}}
\nc{\Ftriv}{\cF_{\mathrm{triv}}}


\nc{\bE}{\underline{E}}
\newcommand{\num}[1]{[#1]}
\newcommand{\dnum}[1]{\num{#1}}%for "dia num"; used in the figures
\DeclareRobustCommand{\stirling}{\genfrac\{\}{0pt}{}}

\nc{\htimes}{\widehat{\otimes}}
\nc{\frakm}{\mathfrak{m}}
%\nc{\frakp}{\mathfrak{p}}
%\usepackage[showframe]{geometry}

\newcommand{\bc}[1]{\mathinner{\left\langle #1 \right\rangle}_{\mkern-2mu *}}
\newcommand{\cbc}[1]{\mathinner{\langle #1 \rangle}^{\mkern-2mu *}}
\newcommand{\bcnothuge}[1]{\mathinner{\Bigl\langle #1 \Bigl\rangle_*}}
\AddToHook{env/Thm/begin}{     \crefalias{equation}{Thm}}
\AddToHook{env/Exa/begin}{     \crefalias{equation}{Exa}}
\AddToHook{env/Prop/begin}{    \crefalias{equation}{Prop}}
\AddToHook{env/Lem/begin}{     \crefalias{equation}{Lem}}
\AddToHook{env/Cor/begin}{     \crefalias{equation}{Cor}}
\AddToHook{env/Conj/begin}{    \crefalias{equation}{Conj}}
\AddToHook{env/Rem/begin}{    \crefalias{equation}{Rem}}
\AddToHook{env/Def/begin}{    \crefalias{equation}{Def}}
\AddToHook{env/Rec/begin}{\crefalias{equation}{Rec}}
\AddToHook{env/War/begin}{\crefalias{equation}{War}}


\title[The tensor triangular geometry of valuation domains]{The tensor triangular geometry \\ of valuation domains}

\author[T.~Barthel]{Tobias Barthel}
\author[D.~Heard]{Drew Heard}
\author[B.~Sanders]{\\Beren Sanders}
\author[C.~Zou]{Changhan Zou}
\date{3 November 2026}

\makeatletter
\patchcmd{\@setaddresses}{\indent}{\noindent}{}{}
\patchcmd{\@setaddresses}{\indent}{\noindent}{}{}
\patchcmd{\@setaddresses}{\indent}{\noindent}{}{}
\patchcmd{\@setaddresses}{\indent}{\noindent}{}{}
\makeatother


\address{Tobias Barthel, Max Planck Institute for Mathematics, Vivatsgasse 7, 53111 Bonn, Germany}
\email{tbarthel@mpim-bonn.mpg.de}
\urladdr{\href{https://sites.google.com/view/tobiasbarthel/home}{https://sites.google.com/view/tobiasbarthel/home}}

\address{Drew Heard, Department of Mathematical Sciences, Norwegian University of Science and Technology, Trondheim}
\email{drew.k.heard@ntnu.no}
\urladdr{\href{https://folk.ntnu.no/drewkh/}{https://folk.ntnu.no/drewkh}}

\address{Beren Sanders, Mathematics Department, UC Santa Cruz, 95064 CA, USA}
\email{beren@ucsc.edu}
\urladdr{\href{https://people.ucsc.edu/~beren/}{https://people.ucsc.edu/$\sim$beren}}

\address{Changhan Zou, Sydney Mathematical Research Institute, The University of Sydney, NSW 2006, Australia}
\email{changhan.zou@sydney.edu.au}
\urladdr{\href{https://maths.usyd.edu.au/ut/people?who=C_Zou}{\nolinkurl{https://maths.usyd.edu.au/ut/people?who=C_Zou}}}

\begin{document}

\begin{abstract}
	We study the tensor triangular geometry of derived categories of valuation domains, obtaining complete ring-theoretic characterizations of their main structural properties. In particular, although valuation domains are almost always non-noetherian, their derived categories always satisfy the local-to-global principle and hence the detection property. We \mbox{establish} precise criteria for homological detection, stratification, and homological \mbox{stratification,} revealing a clear hierarchy among these properties. Notably, we obtain \mbox{examples} of derived categories of commutative rings that are homologically stratified and satisfy the local-to-global principle, but are not stratified. We also determine when the category of modules over a valuation domain has Gabriel dimension, and use this to provide a negative answer to a question of Balchin, G\'omez and Stevenson. Overall, valuation domains provide a rich class of \mbox{non-noetherian} rings whose derived categories can be understood in considerable depth.
\end{abstract}

\vspace*{0.5em}%hack
\maketitle

{
\hypersetup{linkcolor=black}
\tableofcontents
}

\section{Introduction}
The derived category of a commutative ring is a basic example of a tensor \mbox{triangulated} category and it is remarkably well behaved when the ring is \mbox{noetherian.} In this case, its global structure is controlled by theories of support. For example, its localizing subcategories are classified via support by the subsets of the spectrum~\cite{Neeman92a}. In the language of tensor triangular geometry, the derived category is \emph{stratified}. Without the noetherian hypothesis, this can fail dramatically. For example, Neeman~\cite{Neeman00} constructs a ring whose spectrum is a single point, but whose derived category has at least $2^{2^{\aleph_0}}$ Bousfield classes, and hence at least that many localizing subcategories.

Derived categories of non-noetherian rings remain poorly understood in general. A notable exception is the class of absolutely flat rings (the commutative rings of flat dimension zero) which have been comprehensively studied by Stevenson~\cite{Stevenson14a}. The aim of this paper is to carry out an analogous study for valuation domains: the local domains of flat dimension at most one.

The support-theoretic properties of the derived category of a valuation domain~$R$ turn out to be governed by two invariants: the order type of the totally ordered set of prime ideals and the pattern of idempotence among these ideals. A nonzero prime ideal is said to be \emph{branched} if it is not the union of the prime ideals properly contained in it. We say that $R$ is \emph{discrete} if no branched prime ideal is idempotent and \emph{strongly discrete} if no nonzero prime ideal is idempotent. As the name suggests,~$R$ is discrete if and only if its value group is discrete in a certain sense; see~\cref{rem:discrete-iff-discrete}. In particular, a valuation domain of finite Krull dimension is discrete if and only if it is strongly discrete if and only if its value group is $\bbZ^n$ ordered lexicographically. 

Recall that the \emph{homological support} of a complex $X\in\Der(R)$ is defined by
\[
	\Supph(X) = \SETT{\frakp \in \Spec(R)}{k(\frakp) \otimes_R^{\bbL} X \neq 0}
\]
where $k(\frakp)$ denotes the residue field. The \emph{homological detection property} holds for~$\Der(R)$ if $\Supph(X) = \emptyset$ implies $X = 0$. We establish the following characterization of the homological detection property (\cref{thm:h-detection-char}):

\begin{Thm}\label{thm:intro-h-detection-char}
	Let $R$ be a valuation domain. The derived category $\Der(R)$ has the homological detection property if and only if $R$ is discrete.
\end{Thm}

We also determine when the homological support classifies the localizing subcategories of~$\Der(R)$, that is, when the category is \emph{homologically stratified} in the sense of~\cite{BarthelHeardSandersZou26}. The following is~\cref{thm:h-stratification}:

\begin{Thm}\label{thm:intro-h-stratified}
	Let $R$ be a valuation domain. The derived category $\Der(R)$ is homologically stratified if and only if $R$ is discrete and the totally ordered set of prime ideals does not contain a subset order-isomorphic to the rational numbers~$\bbQ$.
\end{Thm}

For derived categories of commutative rings, homological stratification is actually equivalent to the statement that the residue fields generate the category:
\[
	\Der(R) = \Loc\langle k(\frakp) \mid \frakp \in \Spec(R) \rangle.
\]
In other words, homological stratification is equivalent to the homological local-to-global principle. See~\cite[Example~4.10]{BarthelHeardSandersZou26}.


The other support theory used for classification problems in tensor triangular geometry is the Balmer--Favi--Sanders support~\cite{BalmerFavi11,BillySanders17pp} which we recall in~\cref{def:tt-support}. The corresponding notions of detection, local-to-global principle and stratification were developed in~\cite{BarthelHeardSanders23a,Zou25} and they behave very differently from the homological versions. For example, the local-to-global principle always holds (\cref{prop:LGP}):

\begin{Prop}\label{prop:intro-LGP}
	Let $R$ be a valuation domain. The derived category $\Der(R)$ satisfies the local-to-global principle. In particular, it satisfies the detection property.
\end{Prop}

We also determine when the Balmer--Favi--Sanders support classifies the localizing subcategories~(\cref{thm:tt-stratified}):

\begin{Thm}\label{thm:intro-tt-stratified}
	Let $R$ be a valuation domain. The derived category $\Der(R)$ is stratified if and only if $R$ is strongly discrete.
\end{Thm}
 
As a corollary, we recover the following theorem (\cref{cor:tc-char}) which was first proved in~\cite[Theorem~7.2]{BazzoniStovicek17}:

\begin{Thm}[Bazzoni--\Stovicek]\label{thm:intro-telescope}
	Let $R$ be a valuation domain. The telescope conjecture holds for~$\Der(R)$ if and only if $R$ is strongly discrete.
\end{Thm}

We will sometimes refer to stratification as tt-stratification and homological stratification as h-stratification. If the derived category of a commutative ring $R$ is tt-stratified then it is necessarily h-stratified and its spectrum $\Spec(R)$ is necessarily weakly noetherian. The first example of a ring that is h-stratified but not tt-stratified was given in~\cite[Section~3]{BalchinGomezStevenson26pp}. It is a localization of a truncated polynomial ring in infinitely many variables. Our results provide many additional examples among valuation domains. From a broader perspective, these examples demonstrate that homological stratification is the more widely applicable method for classifying localizing ideals.

\begin{figure}[!tb]
	\centering
	\begin{tikzpicture}[region/.style={draw, rounded corners=6pt, thick}, lab/.style={anchor=north west, align=left, font=\small}, ex/.style={anchor=south east, font=\small, text=black}]
	  \draw[region, fill=gray!4]  (0,0)     rectangle (12,7.3);
	  \draw[region, fill=gray!10] (0.3,0.3) rectangle (11.7,5.8);
	  \draw[region, fill=gray!17] (0.6,0.6) rectangle (11.4,4.3);
	  \draw[region, fill=gray!25] (0.9,0.9) rectangle (11.1,2.8);
	  \node[lab] at (0.15,7.15) {\textbf{All valuation domains}\\ $\Der(R)$ satisfies the local-to-global principle};
	  \node[ex]  at (11.85,5.9) {e.g.\ $\Gamma = \bbQ$};
	  \node[lab] at (0.45,5.65) {\textbf{\boldmath$R$ discrete}\\ $\Leftrightarrow$ $\Der(R)$ has the homological detection property};
	  \node[ex]  at (11.55,4.4) {e.g.\ $\Gamma = \prod^{\mathrm{lex}}_{\bbQ}\bbZ$};
	  \node[lab] at (0.75,4.15) {\textbf{\boldmath$R$ discrete and $\Spec(R)$ contains no copy of $\bbQ$}\\ $\Leftrightarrow$ $\Der(R)$ is homologically stratified};
	  \node[ex]  at (11.25,2.9) {e.g.\ $\Gamma = \prod^{\mathrm{lex}}_{\bbN}\bbZ$};
	  \node[lab] at (1.05,2.65) {\textbf{\boldmath$R$ strongly discrete}\\ $\Leftrightarrow$ $\Der(R)$ is stratified\\ $\Leftrightarrow$ the telescope conjecture holds for $\Der(R)$\\ $\Leftrightarrow$ $\Mod(R)$ has Gabriel dimension};
	  \node[ex]  at (10.95,1.0) {e.g.\ $\Gamma = \bbZ^n_{\mathrm{lex}}$};
	\end{tikzpicture}
	\caption{Support-theoretic properties of $\Der(R)$ for a valuation domain $R$ with value group $\Gamma$.}
	\label{fig:summary}
\end{figure}

The following class of examples is illuminating and discussed in more detail in~\cref{sec:value-groups}. For any totally ordered set~$I$, let $R_I$ denote a valuation domain whose value group is the Hahn product~$\prod_I^{\mathrm{lex}} \bbZ$. We will see that this valuation domain is always discrete, and it is strongly discrete if and only if $I\op$ is well-ordered. Moreover, we will see that its ordered set of prime ideals has no copy of $\bbQ$ if and only if the indexing set~$I$ has no copy of $\bbQ$. For example, $R_{\bbN}$ is discrete but not strongly discrete, and of course $\bbN$ does not contain a copy of $\bbQ$. Its derived category is thus h-stratified but not tt-stratified by~\cref{thm:intro-h-stratified} and~\cref{thm:intro-tt-stratified}. This is particularly interesting, since its spectrum coincides with the Balmer spectrum of the stable homotopy category of finite $p$-local spectra (see~\cref{exa:RNN}). Thus, the tensor-triangulated categories $\Der(R_{\bbN})$ and $\SHp$ have the same spectrum of compact objects, but different global properties. The former is h-stratified while the latter does not even have the h-detection property.

Our final contribution (\cref{thm:gabriel-dimension}) is a new proof of the following theorem concerning Gabriel dimension. A proof of this result can also be extracted from the classical ring-theoretic literature; see~\cref{rem:classical}.

\begin{Thm}\label{thm:main-gabriel-dimension}
	Let $R$ be a valuation domain. The category of $R$-modules $\Mod(R)$ has Gabriel dimension if and only if $R$ is strongly discrete.
\end{Thm}

The authors of~\cite{BalchinGomezStevenson26pp} note that they do not know whether $\Mod(R)$ must have Gabriel dimension in order for $\Der(R)$ to be generated by the residue fields (equivalently, to be h-stratified). They also observe that there are few general tools for showing that the residue fields fail to generate the derived category. Our results answer their question negatively. For example, $\Der(R_\bbN)$ is h-stratified but $\Mod(R_\bbN)$ does not have Gabriel dimension.

\medskip
The main results of the paper are summarized in~\cref{fig:summary}.

\medskip
The support-theoretic properties of derived categories of valuation domains have also been studied in the independent and contemporaneous work of \Stovicek~and Tecklenburg \cite{StovicekTecklenburg2026pp}.

\subsection*{Acknowledgments}
TB is supported by the European Research Council (ERC) under Horizon Europe (grant No.~101042990) and would like to thank the Max Planck Institute for its hospitality. 
The first proof of~\cref{thm:one-direction} was obtained via LLMs. The key ingredient was Hausdorff's hierarchy for scattered totally ordered sets, which was unknown to us. The proof presented in this paper is our own and instead uses the notion of Cantor--Bendixson height. LLMs played no role in any other aspect of this work or its results.


\section{Basics on valuation domains}

We begin with a review of the basic properties of valuation domains. A valuable reference for our purposes is~\cite{FuchsSalce01} which can be augmented by the more classical treatment of~\cite{Gilmer92}.

\begin{Def}
	A \emph{valuation domain} is a domain whose set of ideals is totally ordered. Throughout the paper we will be discussing the properties of a valuation domain $R$.
\end{Def}

\begin{Rem}\label{rem:union} 
	The union of any set of prime ideals is itself a prime ideal, since they are totally ordered. It follows that every valuation domain has a unique maximal ideal.
\end{Rem}

\begin{Lem}\label{lem:radical-is-prime}
	A proper ideal is radical if and only if it is prime. In particular, the radical $\sqrt{I}$ of a (nonzero) proper ideal $I$ is a (nonzero) prime ideal.
\end{Lem}

\begin{proof}
	Suppose $I$ is radical and that $ab \in I$. Without loss of generality, we have an inclusion of ideals $(a) \subseteq (b)$. Hence $a =bc$ for some $c \in R$, so that $a^2 = abc \in I$. Hence $a \in I$ since $I$ is radical.
\end{proof}

\begin{Rem}\label{rem:intersection}
	It follows that the intersection of any set of prime ideals is again a prime ideal.
\end{Rem}

\begin{Lem}\label{lem:m-power}
	The intersection $\bigcap_{n \ge 1}\frakm^n$ of powers of the maximal ideal is prime.
\end{Lem}

\begin{proof}
	Let $\frakp \coloneqq \bigcap_{n \ge 1}\frakm^n$ and suppose $ab \in \frakp$. If $a \not\in\frakp$ then $a \not\in \frakm^n$ for some $n \ge 1$. Since ideals are totally ordered, this implies that $\frakm^n \subseteq (a)$. Next observe that for any $k \ge 1$, we have $ab \in \frakm^{n+k} = \frakm^n \frakm^k \subseteq a\frakm^k$ so that $b \in \frakm^k$. Hence $b \in \frakp$.
\end{proof}

\begin{Lem}\label{lem:fg-principal}
	An ideal is finitely generated if and only if it is principal. More precisely, if $I=(a_1,\ldots,a_n)$ then $I=(a_i)$ for some $i=1,\ldots,n$.
\end{Lem}

\begin{proof}
	Since the ideals are totally ordered, there is an $i$ such that $(a_i) \supseteq (a_j)$ for all $1\le j\le n$.
\end{proof}

\begin{Lem}\label{lem:ideal-module-product}
	Let $I$ be an ideal and let $M$ be an $R$-module. If $x \in IM$ then $x = ay$ for some $a\in I$ and $y \in M$.
\end{Lem}

\begin{proof}
	In general, we have $x = \sum_{i=1}^n a_i y_i$ for some $a_i \in I$ and $y_i \in M$. However, since $R$ is a valuation domain, the finitely generated ideal $(a_1,\ldots,a_n)$ must be generated by one of the $a_i$. The claim readily follows.
\end{proof}

\begin{Exa}
	Let $I$ and $J$ be ideals. If $a \in IJ$ then $a=bc$ with $b \in I$ and $c\in J$.
\end{Exa}

\begin{Lem}\label{lem:m-principal}
	If $R$ is not a field then the maximal ideal $\frakm$ is principal if and only if~$\frakm^2 \neq \frakm$.
\end{Lem}

\begin{proof}
	If $\frakm=(a)$ is principal then $\frakm^2=(a^2)$. Hence $\frakm^2=\frakm$ would imply that $a \in (a^2)$ so that $a=a^2 b$ for some $b \in R$. This would imply that $1=ab$ since we are in a domain, which says that $a$ is invertible, which is false. Hence $\frakm^2 \neq \frakm$. Conversely, if $\frakm^2 \neq \frakm$, choose any $a \in \frakm \setminus \frakm^2$. We claim that $\frakm = (a)$. Consider any $b \in \frakm$. If $b \not\in (a)$ then $(a) \subseteq (b)$ since the ideals are totally ordered. Hence $a=bc$ for some $c\in R$. Since $a \not\in \frakm^2$ we must have that $c \not\in \frakm$. This means that $c$ is invertible, so that $b \in (a)$ which contradicts our assumption on $b$.
\end{proof}

\begin{Lem}
	If $S \subset R$ is a multiplicative subset with $0 \not\in S$ then the localization $S^{-1} R$ is again a valuation domain.
\end{Lem}

\begin{proof}
	The lattice of ideals of $S^{-1} R$ is isomorphic to a sublattice of the lattice of ideals of $R$. Hence, the ideals of $S^{-1} R$ are totally ordered. Moreover, $S^{-1} R$ is a domain since $R$ is a domain and $0 \not\in S$.
\end{proof}

\begin{Exa}\label{exa:Rp-vd}
	If $\frakp$ is a prime ideal of $R$ then the localization $R_\frakp$ is again a valuation domain. Similarly, the quotient $R/\frakp$ is a valuation domain.
\end{Exa}

\begin{Lem}\label{lem:dodgy-lemma}
	Let $\frakp$ be a prime ideal. If $a \in \frakp$ and $s \not\in\frakp$ then $a=sb$ for some~${b \in \frakp}$.
\end{Lem}

\begin{proof}
	Note that $s\not\in(a)$. Hence $(a) \subseteq (s)$ since the ideals are totally ordered. Thus $a = bs$ for some $b \in R$. Moreover, $b\in \frakp$ since $a \in \frakp$.
\end{proof}

\begin{Lem}\label{lem:p-idempotent-principal}
	Let $\frakp$ be a nonzero prime ideal of $R$. The following are equivalent:
	\begin{enumerate}
		\item $\frakp^2 \neq \frakp$
		\item The maximal ideal $\frakp R_{\frakp}$ of $R_{\frakp}$ is principal.
	\end{enumerate}
	In this case, $\frakp\Rp = a\Rp$ for some $a \in \frakp$.
\end{Lem}

\begin{proof}
	The equivalence of $(a) \Leftrightarrow (b)$ will follow from \cref{lem:m-principal} and \cref{exa:Rp-vd} once we establish that $\frakp=\frakp^2$ if and only if $\frakp\Rp = \frakp^2\Rp$. Suppose $\frakp\Rp = \frakp^2\Rp$ and let $a \in \frakp$. Then there exist $b,c \in \frakp$ and $s \not\in\frakp$ such that $as=bc$. By \cref{lem:dodgy-lemma}, $c=sd$ for some $d \in \frakp$. Hence $as=bsd$ which implies $a=bd$ since we are in a domain. This establishes that if $\frakp\Rp=\frakp^2\Rp$ then $\frakp = \frakp^2$. The converse is trivial. For the last statement, suppose that $\frakp\Rp$ is principal. Then it is generated by $a/s$ for some $a \in R$ and $s \not\in \frakp$. If $a \not \in\frakp$ then $\frakp\Rp = \frac{a}{s}\Rp=a\Rp=\Rp$ which contradicts the properness of $\frakp\Rp$.
\end{proof}

We now turn to the homological properties of valuation domains.

\begin{Rec} 
	The flat dimension (or weak dimension) of a commutative ring $A$ is the supremum of the minimal lengths of flat resolutions of $A$-modules.
\end{Rec}

\begin{Prop}\label{prop:fd-le-1}
	Let $A$ be a commutative ring. The following are equivalent:
	\begin{enumerate}
		\item The flat dimension of $A$ is $\le 1$.
		\item Every ideal of $A$ is flat.
		\item $\Tor_i^A(M,N) = 0$ for all $i \ge 2$ and all $A$-modules $M$ and $N$.
	\end{enumerate}
\end{Prop}

\begin{proof}
	This is standard; see~\cite[Lemma~4.66]{Lam99} and \cite[Proposition~8.17]{Rotman09}.
\end{proof}

\begin{Def}
	A \emph{Pr\"{u}fer domain} is a domain $A$ with the property that $A_{\frakp}$ is a valuation domain for each prime ideal $\frakp$.
\end{Def}

\begin{Prop}\label{prop:fd-examples} The following statements hold:
	\begin{enumerate}
		\item Commutative rings of flat dimension~0 are precisely the absolutely flat rings.
		\item Domains of flat dimension~0 are precisely the fields.
		\item Domains of flat dimension at most 1 are precisely the Pr\"{u}fer domains.
		\item Local domains of flat dimension at most 1 are precisely the valuation domains.
	\end{enumerate}
\end{Prop}

\begin{proof}
	$(a)$: It is well known that a commutative ring $A$ is absolutely flat (or von Neumann regular) if and only if every $A$-module is flat; see~\cite[Theorem~4.9]{Rotman09}.

	$(b)$: If every $A$-module is flat then $A/aA$ is flat for any $a \in A$. This implies that the injective map $aA \hookrightarrow A$ remains injective after tensoring with $A/aA$. In other words, the zero map $aA/a^2 A \to A/aA$ is injective. This implies that $aA=a^2 A$. If~$A$ is a domain and $a \neq 0$ this implies $A=aA$, so that $a$ is invertible.

	$(c)$: The flat dimension of $A$ is the supremum of the flat dimensions of $A_\frakp$ for each $\frakp \in \Spec(A)$. This follows, for example, from the fact that $\Tor$ commutes with localization \cite[Proposition~7.17]{Rotman09}. Part~$(c)$ thus reduces to part~$(d)$.

	$(d)$: First note that an ideal is always the union of its principal subideals and over a domain every nonzero principal ideal is free of rank 1. The key point is that when $R$ is a valuation domain, the collection of principal subideals forms a filtered system. Hence every ideal is flat, since it is a filtered colimit of flat modules. Thus valuation domains have flat dimension at most 1.

	Conversely, one readily checks that if the principal ideals of a commutative ring are comparable then arbitrary ideals are comparable. We thus claim that a local domain in which every ideal is flat has comparability of principal ideals. To this end, let $a,b\in R$ be two nonzero elements. Consider the ideal $(a,b)$. It is a finitely generated flat $R$-module, and thus free since $R$ is local \cite[Theorem 7.10]{Matsumura89}. Moreover, an ideal cannot have two linearly independent elements. Thus $(a,b)$ is free of rank 1. Therefore $(a,b)=(c)$ for some $c \in R$. We then have $a=cx$, $b=cy$ and $c=za+wb$ for some $x,y,z,w\in R$. It follows that $c=c(xz+yw)$ so that $1=xz+yw$ since $R$ is a domain. In a local ring, the sum of two nonunits is a nonunit. Thus, without loss of generality, $xz$ is a unit, so that $x$ is a unit. Hence $(a)=(c)$ so that $(b) \subseteq (a)$.
\end{proof}

\section{Branched and unbranched prime ideals}\label{sec:branched}

Let $R$ be a valuation domain.

\begin{Rem}\label{rem:union-smaller-p}
	Let $\frakp \subset R$ be a nonzero prime ideal. We can consider the union
	\[
		\frakp'\coloneqq \bigcup_{\frakq \subsetneq \frakp} \frakq
	\]
	of all the prime ideals it strictly contains. This is again a prime ideal, since $R$ is a valuation domain.
\end{Rem}

\begin{Def}
	A nonzero prime ideal $\frakp \subset R$ is said to be \emph{branched} if it is \textbf{not} the union of the prime ideals it strictly contains: $\frakp \neq \frakp'$. A prime ideal which is not branched is said to be \emph{unbranched}. In particular, the zero ideal is unbranched.
\end{Def}

\begin{Lem}\label{lem:branch-equiv}
	Let $R$ be a valuation domain and let $\frakp \subset R$ be a nonzero prime ideal. The following are equivalent:
	\begin{enumerate}
		\item $\frakp = \sqrt{(a)}$ for some $a \in R$.
		\item $\frakp$ is branched.
		\item $\frakp$ has an immediate generalization: there exists $\frakq \in \gen(\frakp)$ with $\frakq\neq \frakp$ such that $\overbar{\{\frakq\}} \cap \gen(\frakp) = \{\frakp,\frakq\}$.
	\end{enumerate}
\end{Lem}

\begin{proof}
	This is established in \cite[Theorem~17.3]{Gilmer92} and is a straightforward exercise. The immediate generalization is of course the $\frakp'$ of~\cref{rem:union-smaller-p}.
\end{proof}

\begin{Lem}\label{lem:radical-unbranched}
	A proper ideal $I$ is prime if its radical $\sqrt{I}$ is unbranched.
\end{Lem}

\begin{proof}
	The statement is trivial for the zero ideal, so assume $I\neq 0$. Suppose $ab\in I$. By~\cref{lem:radical-is-prime}, $\sqrt{I}$ is prime. Hence, without loss of generality, $a \in \sqrt{I}$. Since $\sqrt{I}$ is unbranched, there exists $\frakp \subsetneq \sqrt{I}$ such that $a \in \frakp$. If~$\frakp \not\subseteq I$ then $I \subseteq \frakp$ so that $\sqrt{I} \subseteq \frakp$ which is false. Hence $\frakp \subseteq I$ so $a \in I$.
\end{proof}

\begin{Lem}\label{lem:idempotent-is-prime}
	If a proper ideal is idempotent then it is prime.
\end{Lem}

\begin{proof}
	Suppose $ab \in I$. If $a \not\in I$ then $I \subseteq (a)$ since ideals are totally ordered. Hence $I=I^2 \subseteq aI$ so that $ab \in I \subseteq aI$. Thus, $ab=ac$ for some $c \in I$. Hence $b=c \in I$ since we are in a domain.
\end{proof}

\begin{Lem}\label{lem:unbranched-is-idempotent}
	Every unbranched prime ideal is idempotent.
\end{Lem}

\begin{proof}
	The statement is trivial for the zero ideal, so suppose $a \in \frakp$ is a nonzero element. We claim that $a \in \frakp^2$. Since $\frakp$ is unbranched, it is the union of the prime ideals strictly contained in it. Thus $a \in \frakq$ for some $\frakq \subsetneq \frakp$. Consider any $b \in \frakp$ with~$b \not\in \frakq$. Note that $b \not\in (a)$ since $a \in \frakq$. Thus, $a=bc$ for some $c \in R$. Now $a \in \frakq$, so it follows that $c \in \frakq$. Hence $a\in \frakp\frakq \subseteq \frakp^2$.
\end{proof}

\begin{Def}\label{def:discrete}
	A valuation domain $R$ is said to be
	\begin{enumerate}
		\item \emph{discrete} if no branched prime ideal is idempotent; and
		\item \emph{strongly discrete} if no nonzero prime ideal is idempotent.
	\end{enumerate}
	This terminology will be justified (to some degree) in~\cref{sec:value-groups}. In any case, these notions will be crucial for the remainder of the paper.
\end{Def}

\begin{Prop}\label{prop:discrete-chars} 
	Let $R$ be a valuation domain. The following are equivalent:
	\begin{enumerate}
		\item $R$ is discrete;
		\item The proper idempotent ideals of $R$ are precisely the unbranched prime ideals.
	\end{enumerate}
\end{Prop}

\begin{proof}
	This follows from \cref{lem:idempotent-is-prime} and \cref{lem:unbranched-is-idempotent}.
\end{proof}

\begin{Prop}\label{prop:strongly-discrete-chars}
	Let $R$ be a valuation domain. The following are equivalent:
	\begin{enumerate}
		\item $R$ is strongly discrete;
		\item $R$ is discrete and every nonzero prime ideal is branched.
		\item There are no nonzero proper idempotent ideals.
	\end{enumerate} 
\end{Prop}

\begin{proof}
	This follows from \cref{lem:idempotent-is-prime} and \cref{lem:unbranched-is-idempotent}.
\end{proof}

\begin{Rem}
	We next interpret these notions topologically in terms of the Zariski spectrum. We will use the terminology of~\cite[Section~2]{BarthelHeardSanders23a}.
\end{Rem}

\begin{Lem}\label{lem:thomason-closed}
	Let $S \subseteq \Spec(R)$. The following are equivalent:
	\begin{enumerate}
		\item $S$ is Thomason closed;
		\item $S = V(a)$ for some $a \in R$.
	\end{enumerate}
	Moreover, if $S$ is nonempty and proper, then they are also equivalent to:
	\begin{enumerate}[resume]
		\item $S=\overbar{\{\frakp\}}$ for a branched prime $\frakp$.
	\end{enumerate}
\end{Lem}

\begin{proof}
	For an arbitrary commutative ring, a subset $S\subseteq \Spec(R)$ is Thomason closed if and only if $S=V(I)$ for some finitely generated ideal $I$. In a valuation domain this means $S=V(a)$ for some $a\in R$ by~\cref{lem:fg-principal}. This establishes $(a) \Leftrightarrow (b)$. Next note that if $V(a)$ is nonempty and proper then $a$ is nonzero and a nonunit, hence $\sqrt{(a)} \eqqcolon \frakp$ is a prime ideal by~\cref{lem:radical-is-prime}. Thus $V(a) = V(\sqrt{(a)})=\overbar{\{\frakp\}}$. Moreover,~$\frakp$ is branched by \cref{lem:branch-equiv}. Conversely, if $\frakp$ is a branched prime ideal then it has an immediate generalization $\frakq\subsetneq\frakp$. For any $a \in \frakp \setminus \frakq$, the prime ideal $\sqrt{(a)}$ is contained in $\frakp$ but not in $\frakq$. The immediacy of the inclusion therefore implies that $\sqrt{(a)} = \frakp$. Hence $\overbar{\{\frakp\}} = V(\frakp)=V(a)$.
\end{proof}

\begin{Lem}\label{lem:branched-visible}
	Let $R$ be a valuation domain and let $\frakp \subset R$ be a nonzero prime ideal. The following are equivalent:
	\begin{enumerate}
		\item $\frakp$ is branched;
		\item $\frakp$ is visible;
		\item $\frakp$ is weakly visible.
	\end{enumerate}
\end{Lem}

\begin{proof}
	$(a) \Rightarrow (b)$: By definition, $\frakp$ is visible if $\overbar{\{\frakp\}}$ is Thomason closed. If $\frakp$ is branched then this is the case by \cref{lem:branch-equiv} and \cref{lem:thomason-closed}.

	$(b) \Rightarrow (c)$: This is immediate.

	$(c) \Rightarrow (a)$: If $\frakp$ is weakly visible then $\{\frakp\} = Z \cap \gen(\frakp)$ for some Thomason closed subset $Z$; see~\cite[Remark~2.8]{BarthelHeardSanders23a}. By \cref{lem:branch-equiv} and \cref{lem:thomason-closed}, we have $Z=V(a) = \overbar{\{\frakq\}}$ for the branched prime $\frakq \coloneqq \sqrt{(a)}$. Moreover, $\frakp = \frakq$ since $\{\frakp\} = \overbar{\{\frakq\}} \cap \gen(\frakp)$.
\end{proof}

\begin{Lem}\label{lem:noetherian}
	Let $R$ be a valuation domain with spectrum $\Spec(R)$. The following statements are equivalent:
	\begin{enumerate}
		\item $\Spec(R)$ is noetherian.
		\item $\Spec(R)$ is generically noetherian, i.e., the generalization closure of any point is noetherian. 
		\item $\Spec(R)$ is weakly noetherian.
		\item Every nonzero prime ideal of $R$ is branched.
	\end{enumerate}
\end{Lem}

\begin{proof}
	Condition $(a)$ implies $(b)$ because subspaces of noetherian spaces are noetherian, while~$(b)$ implies $(c)$ by \cite[Lemma 9.9]{BarthelHeardSanders23a}. \Cref{lem:branched-visible} shows that~$(c)$ implies~$(d)$. By \cite[Proposition~7.13]{BalmerFavi11}, $\Spec(R)$ is noetherian if and only if every point is visible. Finally, if $(d)$ holds, then every nonzero prime is visible by \Cref{lem:branched-visible}, and the zero ideal is visible since its closure is all of $\Spec(R)$; hence~$(d)$ implies~$(a)$.
\end{proof}

\begin{Cor}\label{cor:strongly-discrete-noetherian}
	A valuation domain $R$ is strongly discrete if and only if it is discrete and $\Spec(R)$ is (weakly) noetherian. In particular, if $R$ is strongly discrete then $\Spec(R)$ is noetherian.
\end{Cor}

\begin{proof}
	Combine \Cref{prop:strongly-discrete-chars}, $(a)\Leftrightarrow(b)$, with \Cref{lem:noetherian}, $(a)\Leftrightarrow(c)\Leftrightarrow(d)$.
\end{proof}

\section{Value groups and examples}\label{sec:value-groups}

We now briefly review valuations and examples of valuation domains.

\begin{Rec}\label{rec:value-group}
	Let $R$ be a valuation domain with field of fractions~$K$. Its \emph{value group} is the abelian group $\GammaR \coloneqq \KtimesRtimes$ equipped with the total order defined by $\overbar{a} \le \overbar{b}$ iff $b/a \in R$. The canonical projection provides a valuation $v_R\colon K^\times \to \GammaR$ such that $R=\SETT{a \in K^\times}{v_R(a) \ge 0}\cup\{0\}$. Every valuation domain thus arises canonically from a valuation on its field of fractions. Without risk of confusion, we sometimes regard $v_R$ as being defined on the whole of $K$ by setting $v_R(0)=\infty$.
\end{Rec}

\begin{Rem}\label{rem:convex}
	A \emph{convex subgroup} of a totally ordered abelian group $\Gamma$ is a subgroup~$H$ such that if $0 \le b \le c$ in $\Gamma$ and $c\in H$ then $b \in H$. There is an order-reversing bijection between the set of prime ideals of~$R$ and the set of convex subgroups of~$\GammaR$ given by
	\[
		\frakp \mapsto H_\frakp \coloneqq  v_R(R_{\frakp}^{\times}) = \SETT{v_R(a) - v_R(b)}{a,b \in R \setminus \frakp}
	\]
	with inverse $H\mapsto\SETT{a \in R}{v_R(a) > h \text{ for all }h\in H}$; see \cite[Proposition~II.3.4]{FuchsSalce01} and the surrounding discussion. For example, the maximal ideal corresponds to the zero subgroup, and the zero ideal corresponds to the whole group~$\Gamma_R$. The \emph{rank} of a totally ordered abelian group is the supremum of the lengths of chains of convex subgroups. Thus the rank of the value group of a valuation domain $R$ is precisely equal to the Krull dimension of $R$.
\end{Rem}

\begin{Exa}
	The additive groups $\bbZ$, $\bbQ$ and $\bbR$ each have rank 1.
\end{Exa}

\begin{Rem}\label{rem:convex-branched}
	Under the identification between prime ideals and convex subgroups, immediate generalizations $\frakq \subsetneq \frakp$ are identified with immediate inclusions $H_\frakp \subsetneq H_\frakq$. The branched primes therefore correspond to the proper convex subgroups having an immediate successor.
\end{Rem}

\begin{Rem}\label{rem:change-of-group}
	Let $R$ be a valuation domain with value group~$\Gamma$ and let $H_\frakp$ be the convex subgroup associated to a prime ideal $\frakp$. The localization $\Rp$ has value group~$\Gamma/H_\frakp$ and the quotient $R/\frakp$ has value group $H_\frakp$. See~\cite[Lemma~II.3.7]{FuchsSalce01}.
\end{Rem}

\begin{Rec}
	Let $(\Gamma_i)_{i \in I}$ be a family of totally ordered abelian groups whose indexing set $I$ is also totally ordered. Consider the product $\Gamma \coloneqq \prod_{i \in I} \Gamma_i$. The support of an element $g \in \Gamma$ is defined by $\supp(g) \coloneqq \SETT{i \in I}{g_i \neq 0}$. Define 
	\[
		\prod_{i \in I}^{\mathrm{lex}} \Gamma_i = \SETT{ g \in \prod_{i\in I} \Gamma_i}{\supp(g) \text{ is well-ordered}}.
	\]
	This is a subgroup of $\Gamma$ and it becomes a totally ordered abelian group under the following relation: If $g \neq h$ are distinct elements and $i_0$ is the least element of $\supp(h-g)$ then $g \le h$ iff $g_{i_0} \le h_{i_0}$. See \cite[Page~62]{FuchsSalce01} or \cite[Page~167]{Gilmer92}. This is called the \emph{lexicographic product} or the \emph{Hahn product} of the family $(\Gamma_i)_{i \in I}$.
\end{Rec}

\begin{Exa}
	For a finite sequence of totally ordered abelian groups $\Gamma_1,\ldots,\Gamma_n$ indexed by $\{1<2<\cdots<n\}$, this construction gives the usual lexicographic ordering on the product $\Gamma_1 \times \Gamma_2 \times \cdots \times \Gamma_n$. In particular, we can consider the lexicographic product $\bbZ^n_{\mathrm{lex}}$ of~$n$ ordered copies of $\bbZ$. It has rank $n$ as exhibited by
	\[
		0 \subsetneq \{ 0 \}^{n-1} \times \bbZ \subsetneq \{0\}^{n-2} \times \bbZ^2 \subsetneq \cdots \subsetneq \{0\} \times \bbZ^{n-1} \subsetneq \bbZ^n.
	\]
	More generally, we have~\cref{prop:hahn-product-of-archimedean} below.
\end{Exa}

\begin{Rec}
  A \emph{final segment} of a totally ordered set~$I$ is a subset $S \subseteq I$ such that if $i \le j$ and $i \in S$ then $j \in S$. Similarly, we have the notion of an \emph{initial segment} of $I$.
\end{Rec}

\begin{Prop}\label{prop:hahn-product-of-archimedean}
	Let $(\Gamma_i)_{i \in I}$ be a family of rank 1 totally ordered abelian groups indexed by a totally ordered set $I$. The convex subgroups of the Hahn product $\Gamma \coloneqq \prod_{i \in I}^{\mathrm{lex}}\Gamma_i$ are precisely the subgroups
	\[
		H_S=\SETT{g\in\Gamma}{\supp(g)\subseteq S}
	\]
	where $S$ ranges over the final segments of $I$.
\end{Prop}

\begin{proof}
	We first show that the subgroups $H_S$ are convex. Suppose $0 \le g < h$ with $h \in H_S$. If $ g\not\in H_S$ then $\supp(g) \setminus S$ is nonempty. By definition, $\supp(g)$ is well-ordered, so $\supp(g)\setminus S$ has a least element, say $i_0$. Note that $h_{i_0} = 0$ since $i_0 \not\in S$. We claim that $i_0$ is also the least element of $\supp(h-g)$. Indeed, if $j < i_0$ then~$j \not\in S$ since $S$ is a final segment. Thus, $g_j = 0$ by the minimality of $i_0$. Moreover, $h_j = 0$, too. Thus $i_0$ is the least element of $\supp(h-g)$. By the definition of the ordering on the Hahn product, we then have $g_{i_0} \le h_{i_0}$ which is a contradiction. We conclude that $H_S$ is convex.

	Next observe that for any $i \in I$ we have the convex subgroups 
	\[
		H_{\ge i} = \SETT{g \in \Gamma}{g_j = 0 \text{ for } j < i}
			\quad\text{ and }\quad
		H_{> i} = \SETT{g \in \Gamma}{g_j = 0 \text{ for } j \le i}.
	\]
	The quotient $\Gamma \to \Gamma_i$ induces an isomorphism $H_{\ge i}/H_{> i} \simeq \Gamma_i$. Since $\Gamma_i$ is rank~1, there are no convex subgroups sitting between $H_{> i}$ and $H_{\ge i}$. With this in hand, let $H \le \Gamma$ be an arbitrary convex subgroup. Note that the subgroup $H_\emptyset$ corresponding to the empty final segment is the zero subgroup. We may thus assume that $H$ is nonzero. Consider any $0 \neq h \in H$. Let $i = \min\supp(h)$ so that~$h \in H_{\ge i}$. Let $K$ be the convex subgroup of $\Gamma$ generated (as a convex subgroup) by~$h$. We know $K \subseteq H_{\ge i}$ since the latter is convex. Moreover, we can readily check that $H_{> i} \subseteq K$. Thus $K=H_{\ge i}$ since $h \not\in H_{>i}$. It follows that $H_{\ge i} \subseteq H$ since $H$ is convex. Now consider $S\coloneqq \SETT{\min \supp(h)}{0\neq h \in H}$. This is a final segment. Indeed, suppose $i = \min\supp(h)$ with $0\neq h \in H$ as above. For any $j > i$, we can take any nonzero $x \in\Gamma_j$ and let $g \in \Gamma$ be the corresponding element supported at $j$. Then $g \in H_{\ge i} \subseteq H$ and $\min\supp(g)=j$. Hence $j \in S$.

	Finally, note that if $0 \neq h \in H_S$ then $\supp(h) \subseteq S$ by definition. Since $S$ is a final segment, we have $i \coloneqq \min\supp(h) \in S$ and note that $h \in H_{\ge i}$. We then see that $H = \bigcup_{i \in S} H_{\ge i} = H_S$.
\end{proof}

\begin{Rem}\label{rem:convex-of-rank-one-product}
	Note that in the example $\prod_{i \in I}^{\mathrm{lex}} \Gamma_i$ of~\cref{prop:hahn-product-of-archimedean}, the convex subgroups only depend on the indexing set $I$. They have no dependence on the choice of rank~1 groups $\Gamma_i$. Note also that a proper convex subgroup $H_S$ has an immediate successor precisely when the initial segment $I \setminus S$ has a greatest element~$i$. In this case, its successor is $H_{S \cup \{i\}}$. Moreover, observe that $H_{S \cup \{i\}} / H_S \simeq \Gamma_i$.
\end{Rem}

\begin{War}\label{war:discrete}
	There are different definitions of ``discrete'' for totally ordered abelian groups. A common definition is that $\Gamma$ is discrete if it has a least positive element. Of more relevance for our purposes is the following:
\end{War}

\begin{Def}\label{def:discrete-group}
	A totally ordered abelian group~$\Gamma$ is \emph{discrete} if $H/K \simeq \bbZ$ for each immediate inclusion $K \subsetneq H$ in the chain of convex subgroups of $\Gamma$.
\end{Def}

\begin{Exa}
	The lexicographic product $\Gamma \coloneqq \bbQ \times_{\mathrm{lex}} \bbZ$ has a least positive element, namely $(0,1)$. On the other hand, $K= \{0\}\times\bbZ$ is a maximal convex subgroup with~$\Gamma/K \simeq \bbQ$. Thus, $\Gamma$ is not discrete in the sense of \cref{def:discrete-group}.
\end{Exa}

\begin{Exa}\label{exa:hahn-discrete}
	For every totally ordered set $I$, the Hahn product $\prod_I^{\mathrm{lex}} \bbZ$ of $I$ copies of $\bbZ$ is discrete.\footnote{The paragraph after~\cite[Example~2.3]{FuchsSalce01} suggests that the converse holds: that every discrete group is order-isomorphic to a Hahn product of copies of~$\bbZ$. This is false, however. Hill and Mott~\cite[Theorem~3.4]{HillMott73} construct a discrete group (in their terminology, a generalized discrete group) that does not even embed in such a Hahn product.} It has a least positive element only when~$I$ has a greatest element.
\end{Exa}

\begin{Rem}\label{rem:discrete-iff-discrete}
	A valuation domain is discrete in the sense of~\cref{def:discrete} if and only if its value group is discrete in the sense of~\cref{def:discrete-group}. See \cite[Page~205]{Gilmer92}.
\end{Rem}

\begin{Def}\label{def:strongly-discrete-group}
	A totally ordered abelian group $\Gamma$ is \emph{strongly discrete} if it is discrete in the sense of \cref{def:discrete-group} and every proper convex subgroup of $\Gamma$ has an immediate successor in the chain of convex subgroups.
\end{Def}

\begin{Rem}
	A valuation domain $R$ is strongly discrete in the sense of~\cref{def:discrete} if and only if its value group is strongly discrete in the sense of~\cref{def:strongly-discrete-group}. Indeed, $R$ is strongly discrete if and only if $R$ is discrete and every nonzero prime ideal is branched by \cref{prop:strongly-discrete-chars}, so the claim follows from \cref{rem:discrete-iff-discrete} and \cref{rem:convex-branched}.
\end{Rem}

\begin{Rem}\label{rem:hahn-strongly-discrete}
	Let $I$ be a totally ordered set. The Hahn product $\Gamma = \prod_I^{\mathrm{lex}}\bbZ$ is discrete (\cref{exa:hahn-discrete}) and it is strongly discrete if and only if $I\op$ is well-ordered. Indeed, from the definition and~\cref{rem:convex-of-rank-one-product}, we see that $\Gamma$ is strongly discrete if and only if every nonempty initial segment of $I$ has a greatest element. This is equivalent to~$I\op$ being well-ordered, since the initial segment generated by a nonempty subset $S\subseteq I$ has a greatest element precisely when $S$ does.
\end{Rem}

\begin{Exa}\label{exa:lex-N-Z}
	The Hahn product $\prod_{\bbN}^{\mathrm{lex}} \bbZ$ indexed on the countable ordered set~$\bbN$ is discrete but not strongly discrete.
\end{Exa}

\begin{Exa}
	The Hahn product $\prod_{\bbN\op}^{\mathrm{lex}} \bbZ$ indexed on the opposite ordered set $\bbN\op$ is strongly discrete.
\end{Exa}

\begin{Exa}
	A rank $n$ valuation domain is discrete if and only if it is strongly discrete if and only if its value group is $\prod_{\mathrm{lex}}^n \bbZ$. See \cite[Proposition~II.8.4]{FuchsSalce01}.
\end{Exa}

\begin{Rem}
	The following classes of rings coincide:
	\begin{enumerate}
		\item noetherian valuation domains that are not fields;
		\item rank 1 valuation domains with value group $\bbZ$;
		\item rank 1 discrete valuation domains;
		\item rank 1 strongly discrete valuation domains.
	\end{enumerate}
	In particular, we recognize that discrete valuation rings in the usual sense are the rank 1 discrete valuation domains (using the terminology of~\cref{def:discrete}).
\end{Rem}

\begin{Cons}\label{cons:hahn-field}
	Let $\Gamma$ be a totally ordered abelian group and let $K$ be a field. The field of Hahn series $K((t^\Gamma))$ consists of formal power series
	\[
		\sum_{g \in \Gamma} a_g t^g
	\]
	whose support $\SETT{g\in\Gamma}{a_g \neq 0}$ is a well-ordered subset of $\Gamma$. The product is given by convolution using the abelian group structure of $\Gamma$. This convolution is well-defined because the well-ordered hypotheses ensure that for any $g \in \Gamma$ there are only finitely many pairs $(g_1,g_2)$ in the relevant supports such that $g = g_1+g_2$. Additively,~$K((t^\Gamma))$ is the group of functions $\Gamma\to K$ with well-ordered support. There is a canonical surjective valuation $v \colon K((t^\Gamma))^\times \to \Gamma$ defined by $v(f) \coloneqq \min \supp(f)$. We obtain a valuation domain $R = \SETT{ f \in K((t^\Gamma))^\times }{v(f)\ge 0} \cup \{0\}$ whose value group is exactly $\Gamma$. We will refer to~$R$ as the \emph{Hahn valuation domain} associated to~$\Gamma$. Further discussion can be found in \cite[Section~2]{KuhlmannSerra22} or \cite[Section~2]{BalchinTecklenburg25}.
\end{Cons}

\begin{Exa}
	For $\Gamma=\bbZ$, the Hahn series field $K((t^\bbZ))$ is simply the field $K((t))$ of formal Laurent series. Its associated valuation domain is the ring of formal power series $K[[t]]$ which is of course a rank 1 discrete valuation domain.
\end{Exa}

\begin{Exa}
	For a prime $p$ and $\Gamma=\bbZ[1/p]\subseteq \bbQ$, the associated Hahn valuation domain is a rank 1 non-discrete valuation domain. It contains Keller's example from~\cite[\S2]{Keller94b} as a subring with the same value group and residue field, namely the localization of $K[t^{1/p^n} \mid n \ge 0]$ at its augmentation ideal.
\end{Exa}

\begin{Rem}
	Recall that the branched prime ideals correspond to the proper convex subgroups that have an immediate successor. For example, the maximal ideal $\frakm$ is branched if and only if there is a smallest nonzero convex subgroup. On the other hand, one can show that, if $R$ is not a field, the maximal ideal is principal if and only if there is a smallest nonzero convex subgroup \emph{and} that convex subgroup is a copy of $\bbZ$. This, in turn, is the case if and only if the value group $\Gamma$ has a least positive element. This is a two-line exercise using the definitions of~\cref{rec:value-group}. More generally, it then follows from \cref{rem:change-of-group} and \cref{lem:p-idempotent-principal} that a prime ideal~$\frakp$ is not idempotent if and only if $\Gamma/H_\frakp$ has a least positive element. We thus see that the notion of discreteness mentioned in~\cref{war:discrete} is still relevant to the study of valuation domains.
\end{Rem}


\section{Scattered spaces}\label{sec:scattered}

We next discuss the order topology on the set of prime ideals or, equivalently, on the set of convex subgroups of the value group.

\begin{Rec}
	Any totally ordered set~$I$ has an associated \emph{order topology} with subbase given by the open rays $\SETT{x}{x < a}$ and $\SETT{x}{x > a}$ for all $a \in I$. This topology is always Hausdorff~\cite[Theorem~1]{Erne80} and if $I$ is complete, meaning that every subset has a supremum and an infimum, then it is quasi-compact~\cite[Theorem~9]{Frink42}. Note that the order topology of a totally ordered set coincides with its interval topology. This is a simple exercise, or see \cite[Theorem~3]{Frink42}. Recall that the prime ideals of a valuation domain form a complete totally ordered set by \cref{rem:union} and \cref{rem:intersection}.
\end{Rec}

\begin{Lem}\label{lem:constructible}
	Let $R$ be a valuation domain. The order topology on $\Spec(R)$ coincides with the constructible topology on $\Spec(R)$.
\end{Lem}

\begin{proof}
	This is~\cite[Theorem~1.6.4]{DickmannSchwartzTressl19}.
\end{proof}

\begin{Rem}
	A totally ordered set is said to be \emph{dense} if each pair of distinct elements has an element lying between them. For example, the rationals $\bbQ$ and the real numbers $\bbR$ are dense totally ordered sets, while the integers $\bbZ$ are not. A totally ordered set is said to be \emph{scattered} if it contains no densely ordered subset with more than one element. This is equivalent\footnote{The nontrivial direction of the equivalence utilizes Cantor's Theorem~\cite[Theorem II.2.5]{Rosenstein82} which characterizes $\bbQ$ as the unique nonempty countable dense totally ordered set which has no least element and no greatest element. Suppose $A$ is a totally ordered set which is not scattered. By definition it has a dense suborder $D$ which has at least two distinct elements $x < y$. The open interval $(x,y)$ in $D$ is still dense and has no least or greatest element, and we can inductively construct a countable dense suborder as follows. Choose two elements $x_0<x_1$ in $(x,y)$ and set $D_0 \coloneqq \{x_0 < x_1\}$. Given $D_n = \{d_0 < d_1 < \cdots < d_k\}$ for some $k \geq 1$, we construct $D_{n+1}$ by inserting a new element of $(x,y)$ below $d_0$, a new element above $d_k$, and a new element between $d_{i-1}$ and $d_i$ for each $i=1,\ldots,k$. Then $\bigcup_{n=0}^\infty D_n$ is a countable dense suborder which has no minimum element and no maximum element. Hence it is order-isomorphic to the rationals by Cantor's Theorem.} to having no suborder isomorphic to the rationals $\bbQ$.
\end{Rem}

\begin{Rem}
	It follows from the definition that a totally ordered set~$I$ is scattered if and only if its opposite order~$I\op$ is scattered.
\end{Rem}

\begin{Rem}\label{rem:initial-scattered}
	Let $\cat I$ denote the set of initial segments of a totally ordered set~$I$. It is totally ordered by inclusion and we have an embedding $I \hookrightarrow \cat I$ given by $i \mapsto I_{\le i} \coloneqq \SETT{j \in I}{j \le i}$. Using this embedding, it is a straightforward exercise to show that $I$ is scattered if and only if $\cat I$ is scattered.
\end{Rem}

\begin{Rec}\label{def:scattered}
	Let $S$ be a subset of a topological space~$X$. A point $x \in S$ is an \emph{isolated point} of $S$ if $\{x\} = U \cap S$ for some open subset $U \subseteq X$. In other words, $x \in S$ is isolated if $\{x\}$ is open in the subspace topology on $S$.

\end{Rec}

\begin{Def}
	A topological space is said to be \emph{scattered} if every nonempty closed subset has an isolated point.
\end{Def}

\begin{Rem}\label{rem:scattered-all-subsets}
	A topological space $X$ is scattered if and only if every nonempty subset $S \subseteq X$ has an isolated point. One direction is clear. For the other, if $x$ is an isolated point of $\overbar{S}$, then $x \in S$ and $x$ is an isolated point of $S$.
\end{Rem}

\begin{Lem}\label{lem:totally-scattered}
	Let $R$ be a valuation domain. The totally ordered set of prime ideals is scattered if and only if the order topology on $\Spec(R)$ is scattered.
\end{Lem}

\begin{proof}
	It is generally true that the order topology on a scattered totally ordered set is scattered in the topological sense. The converse does not hold in general, but it does hold for complete totally ordered sets, such as $\Spec(R)$. This is a classical result due to Cantor and Hausdorff. See~\cite[Section~2]{PouzetKaddourZaguia14}, for example.
\end{proof}

\begin{Rec}\label{def:CBtodo}
	The \emph{Cantor--Bendixson derivative} of a topological space~$X$ is the subspace $\delta X$ of non-isolated points of $X$. More generally,  the Cantor--Bendixson derivatives of $X$ are defined by transfinite recursion, as follows. We set $\delta^0 X \coloneqq X$ and for every ordinal~$\alpha$, we define $\delta^{\alpha+1} X \coloneqq \delta(\delta^\alpha X)$. Finally, for every nonzero limit ordinal $\lambda$, we set $\delta^\lambda X \coloneqq \bigcap_{\alpha < \lambda} \delta^\alpha X$.
\end{Rec} 

\begin{Lem}\label{lem:scattered-CB}
	A topological space $X$ is scattered if and only if $\delta^\alpha X = \emptyset$ for some ordinal $\alpha$.
\end{Lem}

\begin{proof}
	Suppose $X$ is scattered. Note that if $\delta^\alpha X\neq \emptyset$ then it has an isolated point by \cref{rem:scattered-all-subsets} and hence $\delta^{\alpha+1}X \subsetneq \delta^\alpha X$. Therefore, we can't have $\delta^\alpha X \neq \emptyset$ for all ordinals $\alpha$, because we would then have a proper class of distinct elements in $X$.

	Conversely, suppose $\delta^\alpha X = \emptyset$ for some $\alpha$ and let $S\subseteq X$ be a nonempty subset. Let $\beta$ be the least ordinal such that $\beta < \alpha$ and $S \not\subseteq \delta^{\beta+1} X$. The existence of an ordinal satisfying these two properties is a straightforward consequence of the fact that $S \not\subseteq \delta^\alpha X$. The minimality of $\beta$ readily implies that $S \subseteq \delta^\beta X$. We can then choose an $x \in S$ with $x \not\in \delta^{\beta +1}X$. This point is isolated in $\delta^\beta X$ and hence is isolated in $S$. This establishes that $X$ is scattered.
\end{proof}

\begin{Def}
	The \emph{Cantor--Bendixson height} of a scattered space $X$ is
	\[
		\CBheight(X) \coloneqq \min\{ \alpha \mid \delta^\alpha X = \emptyset \}.
	\]
\end{Def}

\begin{Lem}
	If $A$ is a closed subspace of a topological space $X$ then $A \setminus \Iso(A)$ is closed in $X$. Here $\Iso(A)$ denotes the isolated points of the space $A$.
\end{Lem}

\begin{proof}
	For each $a\in \Iso(A)$, the singleton $\{a\}$ is open in the subspace $A$. Thus, $\Iso(A)$ is open in $A$. Hence, $A \setminus \Iso(A)$ is closed in $A$ and therefore closed in $X$.
\end{proof}

\begin{Lem}\label{lem:delta-closed}
	For every ordinal $\alpha$, the Cantor--Bendixson derivative $\delta^\alpha X$ is a closed subset of $X$.
\end{Lem}

\begin{proof}
	We use transfinite induction on $\alpha$. The base case $\alpha=0$ is clear. Suppose $\delta^\alpha X$ is closed. Then $\delta^{\alpha+1} X=\delta^\alpha X\setminus\Iso(\delta^\alpha X)$ is closed by the previous lemma. For a limit ordinal $\lambda$, we have $\delta^\lambda X= \bigcap_{\beta < \lambda} \delta^\beta X$ which is an intersection of closed subsets.
\end{proof}

\begin{Prop}\label{prop:CB-successor}
	If $X$ is a nonempty scattered  quasi-compact space then the Cantor--Bendixson height is a successor ordinal, say $\CBheight(X)=\alpha+1$, and $\delta^\alpha X$ is finite.
\end{Prop}

\begin{proof}
 	Let $\gamma$ denote the Cantor--Bendixson height of $X$. Suppose for contradiction that $\gamma$ is a limit ordinal. Then
    \[
		\emptyset=\delta^\gamma X= \bigcap_{\beta<\gamma}\delta^\beta X.
    \]
	By the minimality of $\gamma$ and \cref{lem:delta-closed}, $\{\delta^\beta X\}_{\beta<\gamma}$ is a family of closed subsets with the finite intersection property. Quasi-compactness of $X$ then implies $\delta^\gamma X\neq\emptyset$, a contradiction. Hence, $\gamma$ is a successor ordinal. Put $\gamma=\alpha+1$. It remains to show that $\delta^\alpha X$ is finite. This follows from the fact that the discrete space $\delta^\alpha X=\Iso(\delta^\alpha X)$ is quasi-compact, being a closed subspace of the quasi-compact space $X$.
\end{proof}

\begin{Lem}\label{lem:delta-I}
	Let $X$ be a topological space. For any subspace $I \subseteq X$, we have $\delta^\alpha I \subseteq I \cap \delta^\alpha X$ for any ordinal $\alpha$.
\end{Lem}

\begin{proof}
    We use transfinite induction on $\alpha$. The $\alpha=0$ case is trivial. Suppose $\delta^\alpha I \subseteq I \cap \delta^\alpha X$. Let $x\in\delta^{\alpha+1}I$. Then $x\in\delta^\alpha I\subseteq I$. If $x\not\in\delta^{\alpha+1}X$ then $x$ is isolated in $\delta^\alpha X$ and therefore also isolated in $\delta^\alpha I$, contradicting $x\in\delta^{\alpha+1}I$. Therefore $x\in\delta^{\alpha+1}X$ and hence $\delta^{\alpha+1}I\subseteq I \cap \delta^{\alpha+1}X$. If $\lambda$ is a limit ordinal, then
    \[
		\delta^\lambda I= \bigcap_{\beta<\lambda}\delta^\beta I \subseteq \bigcap_{\beta<\lambda} (I\cap \delta^\beta X) =I\cap\delta^\lambda X.
    \]
    This completes the induction.
\end{proof}

\begin{Prop}\label{prop:drop-degree}
	Let $X$ be a nonempty scattered quasi-compact space with $\CBheight(X) = \alpha + 1$. If $I \subseteq X$ is a nonempty subspace such that $I \cap \delta^\alpha X = \emptyset$ then $\delta^\alpha I = \emptyset$. Consequently, $\CBheight(I) < \CBheight(X)$.
\end{Prop}

\begin{proof}
    This follows from \cref{prop:CB-successor} and \cref{lem:delta-I}.
\end{proof}

\section{Homological detection}

We next characterize when the derived category of a valuation domain has the homological detection property.

\begin{Rec}\label{rec:h-support}
	Let $A$ be a commutative ring and let~$\Der(A)$ denote its derived category. The \emph{homological support} of a complex~$X \in \Der(A)$ is
	\[
		\Supph(X) = \SETT{\frakp \in \Spec(A)}{k(\frakp) \otimes_A^{\bbL} X \neq 0 \text{ in } \Der(k(\frakp))}
	\]
	where $k(\frakp) \coloneqq A_{\frakp}/\frakp A_{\frakp}$ is the residue field at~$\frakp$. We say that the \emph{homological detection property} (or \emph{h-detection property}) holds for $\Der(A)$ if $\Supph(X) = \emptyset$ implies $X = 0$.
\end{Rec}

\begin{Rem}\label{rem:localization-of-hdp}
	The homological detection property is studied in the broader generality of rigidly-compactly generated tensor-triangulated categories in~\cite{BarthelHeardSandersZou26}. It is formal that if $\Der(A)$ satisfies the homological detection property then so does $\Der(A_{\frakp})$ for any prime ideal $\frakp \subset A$.
\end{Rem}

\begin{Lem}\label{lem:annihilator-prime}
	Let $R$ be a valuation domain with fraction field $K$. Let $M$ be an $R$-module such that $M \otimes_R K =0$. For any nonzero $x \in M$, the radical of its annihilator 
	\[
		\frakp_x \coloneqq \sqrt{\Ann_R(x)}
	\]
	is a nonzero prime ideal. Moreover, the element $x$ remains nonzero in $M_{\frakp_x}$.
\end{Lem}

\begin{proof}
  Since $M \otimes_R K = 0$ by hypothesis, every element of $M$ is torsion. Hence the annihilator $\Ann_R(x)$ of any element is a nonzero ideal, and it is proper provided $x \neq 0$. In a valuation domain, the radical of any proper nonzero ideal is a nonzero prime ideal (\cref{lem:radical-is-prime}). Finally, if $x/1 =0$ in the localization $M_{\frakp_x}$ then there exists some $s \not\in \frakp_x$ such that $sx=0$ in $M$. This would mean that $s \in \Ann_R(x) \subseteq \frakp_x$ which is false.
\end{proof}

\begin{Prop}\label{prop:detection-for-modules}
  Let $R$ be a valuation domain which is discrete and let $M$ be an $R$-module. If $M \otimes_R^{\bbL}  k(\frakp) = 0$ for all prime ideals $\frakp$, including the zero ideal, then~$M=0$.
\end{Prop}

\begin{proof}
    Suppose for a contradiction that $M \neq 0$. Since $K = k(0)$ is flat, the hypothesis at $\frakp = 0$ gives $M \otimes_R K = M \otimes_R^{\bbL} k(0) = 0$. Hence, by~\cref{lem:annihilator-prime}, each nonzero element $x \in M$ has an associated nonzero prime ideal $\frakp_x = \sqrt{\Ann_R(x)}$. There are two cases to consider.

	\smallskip
	Case (1): Suppose there exists a nonzero $x \in M$ such that $\frakp\coloneqq\frakp_x$ is not idempotent. Hence, by \cref{lem:p-idempotent-principal}, the maximal ideal $\frakp R_{\frakp}$ of $R_{\frakp}$ is principal: there is an element $a \in \frakp$ such that $a R_{\frakp} = \frakp R_{\frakp}$. By \cref{lem:annihilator-prime}, the $R_{\frakp}$-module $M_{\frakp}$ is nonzero; more precisely, $x/1$ is a nonzero element. Now consider the residue field $k(\frakp) = R_{\frakp}/\frakp R_{\frakp} = R_{\frakp}/aR_{\frakp}$. Observe that 
	\[
		0 \to R_{\frakp} \xrightarrow{a} R_{\frakp} \to k(\frakp) \to 0
	\]
	is a free resolution of $k(\frakp)$ as an $R_{\frakp}$-module. Using the fact that Tor commutes with localization, we thus have
	\[
		\Tor_1^R(M,k(\frakp)) = \Tor_1^{R_{\frakp}}(M_{\frakp},k(\frakp)) = \ker(M_{\frakp} \xrightarrow{a} M_{\frakp}).
	\]
	By hypothesis, this Tor group vanishes. Hence
	\begin{equation}\label{eq:inj}
		a\colon M_{\frakp} \to M_{\frakp}
	\end{equation}
	is injective. But recall that $a \in \frakp = \sqrt{\Ann_R(x)}$. Thus $a^n x = 0$ for some $n \ge 1$. Hence $a^n (x/1) = 0$ in $M_{\frakp}$. Since \eqref{eq:inj} is injective, it follows that $x/1=0$ which is a contradiction.
  
	Case (2): Suppose that $\frakp_x$ is idempotent for every nonzero $x \in M$. By our hypothesis on $R$, this idempotent ideal $\frakp_x = \sqrt{\Ann_R(x)}$ must be unbranched. Hence $\Ann_R(x)$ is itself a prime ideal by~\cref{lem:radical-unbranched}. Thus, we are in the situation where the annihilator of every nonzero element of $M$ is prime. We emphasize that this is true for every nonzero element. Nevertheless, choose any nonzero $x \in M$ and set $\frakp \coloneqq \frakp_x = \Ann_R(x)$. We will prove that $M \otimes_R k(\frakp)\neq 0$, yielding a contradiction. As before, $x/1$ is nonzero in the $R_{\frakp}$-module $M_{\frakp}$. We claim that the image of $x/1$ in
	\[
		M_\frakp/\frakp M_\frakp \simeq M\otimes_R k(\frakp)
	\]
	is nonzero. Clearing denominators gives $sx \in \frakp M$ for some $s \notin \frakp$. By \cref{lem:ideal-module-product}, we can therefore write
	\begin{equation}
		sx = ay
	\end{equation}
	for some $a \in \frakp$ and $y \in M$. Note that $ay\neq 0$ since otherwise we would have $sx=0$ but $s \not\in\frakp=\Ann_R(x)$. In particular, $y \neq 0$. Recall that we are in the situation where the annihilator of every nonzero element of $M$ is prime. In particular, $\frakq \coloneqq \frakp_{y} = \Ann_R(y)$ is a nonzero prime ideal. Note that $a \not\in \frakq$ since $ay\neq 0$. On the other hand, since $a \in \frakp$, we have $ax=0$. Hence $asx=0$ so that $a^2y=0$. Hence $a^2 \in \frakq$. Since $\frakq$ is prime, this implies $a \in \frakq$, which is a contradiction.
\end{proof}

\begin{Prop}\label{prop:detection-for-complexes}
	Let $R$ be a valuation domain. If $R$ is discrete then $\Der(R)$ has the h-detection property.
\end{Prop}

\begin{proof}
	Let $X \in \Der(R)$ be a complex. Suppose that $X \otimes_R^{\bbL} k(\frakp) = 0$ for all prime ideals $\frakp$. We have the Tor spectral sequence 
	\[
		E_{pq}^2 = \Tor^R_p(H_q(X),k(\frakp)) \Longrightarrow H_{p+q}(X \otimes_R^{\bbL} k(\frakp)).
	\]
	Since $R$ has flat dimension at most 1 (\cref{prop:fd-examples}), the $E^2$-page is concentrated in the columns $p=0$ and $p=1$. Hence all differentials $d^r$ ($r \ge 2$) vanish, and convergence yields short exact sequences
	\[
		0 \to H_i(X) \otimes_R k(\frakp) \to H_i\big(X \otimes_R^{\bbL} k(\frakp)\big) \to \Tor_1^R\big(H_{i-1}(X),k(\frakp)\big) \to 0.
	\]
	By hypothesis, the middle term vanishes. Hence
	\[
		H_i(X) \otimes_R k(\frakp) = 0 \quad\text{ and }\quad \Tor_1^R(H_{i}(X),k(\frakp)) = 0
	\]
	for all $i$. Since the higher Tor terms vanish, this is equivalent to saying that $H_i(X) \otimes_R^{\bbL} k(\frakp) = 0$. Thus $H_i(X) = 0$ by~\cref{prop:detection-for-modules}. This is true for all $i$ so~$X=0$ in $\Der(R)$. 
\end{proof}

\begin{Prop}\label{prop:h-detection-implies-discrete}
	Let $R$ be a valuation domain. If the derived category $\Der(R)$ has the h-detection property then $R$ is discrete.
\end{Prop}

\begin{proof}
	We prove the contrapositive. Suppose that $R$ is not discrete. This means that there exists a branched prime $\frakp$ with $\frakp^2=\frakp$. We claim that $\Der(R)$ does \emph{not} have the h-detection property. If it \emph{did} have the h-detection property then so would~$\Der(R_{\frakp})$ by~\cref{rem:localization-of-hdp}. Thus it suffices to prove that $\Der(R_{\frakp})$ does not have the h-detection property. Moreover, the maximal ideal of the localization $R_{\frakp}$ is still branched and idempotent by~\cref{lem:p-idempotent-principal} and~\cref{lem:branch-equiv}. (Note that $R_{\frakp}$ is not a field since $\frakp \neq 0$.) It follows that we may assume without loss of generality that $\frakp$ is the maximal ideal $\frakm$.

	That is, it suffices to prove that if $R$ is a valuation domain whose maximal $\frakm$ ideal is branched and idempotent then the h-detection property does not hold. Since $\frakm$ is branched, there exists an element $a \in \frakm$ which is not contained in any nonmaximal prime ideal. Moreover, since $\frakm$ is idempotent, it is not a principal ideal (\cref{lem:m-principal}). Thus $aR \subsetneq \frakm$. In particular, we see that $a\frakm \subsetneq \frakm$ is a proper subideal. We claim that the nonzero $R$-module $M\coloneqq \frakm / a\frakm$ has empty homological support. Since ideals are flat, we have a resolution
	\[
		0 \to a\frakm \to \frakm \to \frakm/a\frakm \to 0
	\]
	by flat modules. Let $\frakq$ be any prime ideal. We claim that $\frakm/a\frakm \otimes_R^{\bbL} k(\frakq)=0$; in other words, that the complex
	\[
		0 \to a\frakm\otimes_R k(\frakq) \to \frakm\otimes_R k(\frakq) \to 0
	\]
	is acyclic.

	For the maximal ideal $\frakq=\frakm$, we have $a\frakm \otimes_R R/\frakm =0$ and $\frakm \otimes_R R/\frakm =0$ since~$\frakm$ is idempotent. Hence the above complex vanishes.

	On the other hand, for any prime ideal $\frakq \neq \frakm$, we have that $a \not\in\frakq$ by construction. It follows that the map $a\frakm \to \frakm$ becomes an isomorphism when we tensor with $R_{\frakq}$ (and hence when we tensor with $k(\frakq)$). Indeed, the map remains injective since $R_{\frakq}$ is flat. It is surjective since $a$ is one of the elements which gets inverted.

	In summary, $\frakm/a\frakm \otimes_R^{\bbL} k(\frakq)=0$ for all prime ideals $\frakq$.
\end{proof}

\begin{Thm}\label{thm:h-detection-char}
	Let $R$ be a valuation domain. The derived category $\Der(R)$ has the homological detection property if and only if $R$ is discrete (meaning that none of the branched primes of $R$ are idempotent).
\end{Thm}

\begin{proof}
	This is established by \cref{prop:detection-for-complexes} and \cref{prop:h-detection-implies-discrete}.
\end{proof}

\section{Balmer--Favi idempotents}

We next consider the Balmer--Favi theory of support~\cite{BalmerFavi11} and the associated theory of stratification developed in~\cite{BarthelHeardSanders23a} and \cite{BillySanders17pp,Zou25}.

\begin{Rec}\label{rec:balmer-favi}
	Let $A$ be a commutative ring. Associated to any Thomason subset~$Y\subseteq\Spec(A)$ is a smashing localization of $\Der(A)$ whose kernel is generated by the perfect complexes supported on $Y$. We write $e_Y \to A \to f_Y \to \Sigma e_Y$ for the associated idempotent triangle. A subset ${W\subseteq\Spec(A)}$ is said to be \emph{weakly visible} if $W=Y_1\cap Y_2^{\complement}$ for Thomason subsets $Y_1$ and $Y_2$. The idempotent object $\gW \coloneqq e_{Y_1} \otimes_A^{\bbL} f_{Y_2}$ does not depend on the choice of $Y_1$ and $Y_2$ by \cite[Remark~7.6]{BalmerFavi11}. We simply write $g_\frakp \coloneqq g_{\{\frakp\}}$ when the singleton $\{\frakp\}$ is weakly visible. See~\cite[Section~4]{BarthelCastellanaHeardSanders23app} for further discussion.
\end{Rec}

\begin{Def}\label{def:tt-support}
    The \emph{Balmer--Favi--Sanders support} is defined for $X \in \Der(A)$ by 
	\[
		\Supp(X) \coloneqq \SETT{\frakp \in \Spec(A)}{g_W \otimes_A^{\bbL} X \neq 0 \text{ for every weakly visible } W \ni \frakp}.
	\]
	We say that the \emph{detection property} holds if $\Supp(X) = \emptyset \implies X = 0$. We also say that the \emph{local-to-global principle} holds if
	\[
		A \in \Loco{g_{W_i} \mid i \in I}
	\]
	for every cover of $\Spec(A)$ by weakly visible subsets $W_i$; see~\cite[Definition~7.1 and Remark~7.2]{Zou25}. Finally, we say that $\Der(A)$ is \emph{stratified} (or \emph{tt-stratified}) if the map 
	\[
		\{\text{localizing ideals of }\Der(A) \} \xrightarrow{\Supp} \{\text{localizing closed subsets of }\Spec(A)\}
	\]
	is a bijection.
\end{Def}

\begin{Rem}
	If $\Spec(A)$ is weakly noetherian then every subset of $\Spec(A)$ is localizing closed\footnote{See~\cite[Remark~2.7]{Zou25} for the precise definition.} and these definitions reduce to those studied in~\cite{BarthelHeardSanders23a}. In fact, if~$\Der(A)$ is stratified then $\Spec(A)$ must be weakly noetherian by~\cite[Theorem~A]{Zou25}. Thus, the extra generality is not required for stratification. However, the detection property and the local-to-global principle can hold more generally, as we will see below.
\end{Rem}

For the next result, we need the following definition. 

\begin{Rec}\label{rec:weakly-scattered}
	Let $S$ be a subset of a topological space $X$. A point $x \in S$ is a \emph{weakly isolated point} of $S$ if $\{x\} \subseteq U \cap S \subseteq \overbar{\{x\}}$ for some open subset $U \subseteq X$. The space $X$ is \emph{weakly scattered} if every nonempty closed subset of $X$ has a weakly isolated point. Since isolated points are weakly isolated, scattered spaces (\cref{def:scattered}) are weakly scattered. A spectral space $X$ is \emph{Hochster (weakly) scattered} if its Hochster dual $X^\vee$ is (weakly) scattered. Recall that the open subsets of $X^\vee$ are the Thomason subsets of $X$, so the closure of a point $x$ in $X^\vee$ is the set~$\gen(x)$ of generalizations of~$x$.
\end{Rec}

\begin{Prop}\label{prop:LGP}
	For every valuation domain $R$, the derived category $\Der(R)$ satisfies the local-to-global principle. In particular, it satisfies the detection property.
\end{Prop}

\begin{proof}
	By~\cite[Theorem~7.6]{Zou25}, it suffices to show that $\Spec(R)$ is Hochster weakly scattered. To this end, let $Y$ be any proper Thomason subset of $\Spec(R)$. We need to find a prime ideal $\frakp$ that is not contained in $Y$ and a Thomason subset $U$ such that $\frakp \in U \cap Y^\complement \subseteq \gen(\frakp)$. Consider the union of all prime ideals not in $Y$:
	\[
		\frakp \coloneqq \bigcup_{\frakq \not\in Y} \frakq.
	\]
	Since $R$ is a valuation domain, $\frakp$ is indeed a prime ideal. Note that $Y^\complement \subseteq \gen(\frakp)$ by the definition of $\frakp$. We claim that $\frakp \not\in Y$. Indeed, if $\frakp \in Y$ then we would have $\frakp \in V(a) \subseteq Y$ for some element $a \in R$ by~\Cref{lem:thomason-closed}. The definition of~$\frakp$ then implies that there exists some $\frakq \not\in Y$ such that $\frakq \in V(a) \subseteq Y$ which is a contradiction. Hence, we have $\frakp \in \Spec(R) \cap Y^\complement = Y^\complement \subseteq \gen(\frakp)$. This proves that $\Spec(R)$ is Hochster weakly scattered.
\end{proof}

\begin{Rem}
	The detection property does not hold for all commutative rings; see our upcoming work \cite{BarthelHeardSandersZou26pp}.
\end{Rem}

\begin{Lem}\label{lem:for-one-for-all}
	Let $R$ be a valuation domain which is not a field. Suppose that the maximal ideal is principal, say $\frakm = (a)$. Then we have $R[1/a]=\Rp$ where $\frakp \subsetneq \frakm$ is the immediate generalization of the maximal ideal.
\end{Lem}

\begin{proof}
	Note that the (nonzero) $\frakm$ is branched by \cref{lem:m-principal} and \cref{lem:unbranched-is-idempotent}. Hence it does indeed have an immediate generalization~$\frakp$. Note that $a \notin \frakp$, since otherwise $\frakm = (a) \subseteq \frakp$. The intersection $\frakq \coloneqq \bigcap_{n \ge 1}\frakm^n$ is a prime ideal by~\cref{lem:m-power} and it does not equal $\frakm$ since that would imply $\frakm^2=\frakm$. It follows that $\frakq \subseteq \frakp$ since $\frakp$ is an immediate generalization of $\frakm$ and ideals are totally ordered. Now consider~$x\in\frakp$. If $x\not\in\frakq$ then we can consider the largest $n \ge 1$ such that $x\in \frakm^n$. We have $x=a^n y$ for some $y\in R$ but since $\frakp$ is prime and $a \not \in \frakp$, we must have $y \in \frakp \subseteq \frakm$. However, this implies that $x \in \frakm^{n+1}$ which is false. We have thus established that~$\frakp=\frakq$. Now if $s$ is an element which is not contained in $\frakp$ then $s \not\in \frakm^n$ for some~$n \ge 1$. Hence~$\frakm^n \subseteq (s)$. Thus $a^n =s b$ for some $b \in R$. We then see that in the localization~$R[1/a]$, the element $b/a^n$ provides an inverse for $s$. In other words, every element $s \notin \frakp$ becomes invertible in $R[1/a]$, so that $\Rp \subseteq R[1/a]$. Conversely, $R[1/a] \subseteq \Rp$ since $a \notin \frakp$.
\end{proof}

\begin{Prop}\label{prop:gp-kp}
	Let $R$ be a valuation domain and let $\frakp$ be a branched prime ideal. Then
	\begin{equation}\label{eq:gp-tensor-kp}
		g_\frakp \otimes_R^{\bbL} k(\frakp) \simeq k(\frakp).
	\end{equation}
	If moreover $\frakp$ is not idempotent, then
	\begin{equation}\label{eq:gp-kp}
		\Loco{g_\frakp} = \Loco{k(\frakp)}.
	\end{equation}
\end{Prop}

\begin{proof}
	First note that~$\frakp$ is visible by~\cref{lem:branched-visible}, so $g_\frakp$ is defined. In general, if $\{\frakp\} = Y_1 \cap Y_2^\complement$ for Thomason subsets $Y_1, Y_2 \subseteq \Spec(R)$ then tensoring the idempotent triangle $e_{Y_1} \to \unit \to f_{Y_1} \to \Sigma e_{Y_1}$ with $f_{Y_2}$ gives an exact triangle
	\begin{equation}\label{gp-triangle}
		g_\frakp \to f_{Y_2} \to f_{Y_1 \cup Y_2} \to \Sigma g_\frakp.
	\end{equation}
	We take $Y_1 = \overbar{\{\frakp\}}$, which is Thomason closed since $\frakp$ is visible, and $Y_2 = \gen(\frakp)^\complement$, which is a Thomason subset, being the union of the $V(a)$ with $a \notin \frakp$. By~\cref{lem:branch-equiv},~$\frakp$ has an immediate generalization $\frakq \subsetneq \frakp$, and since the prime ideals are totally ordered, $Y_1 \cup Y_2 = \overbar{\{\frakp\}} = \gen(\frakq)^\complement$. Since the finite localization away from $\gen(x)^\complement$ is localization at $x$, the triangle~\eqref{gp-triangle} becomes $g_\frakp \to \Rp \to \Rq \to \Sigma g_\frakp$, where $\Rp \to \Rq$ is the canonical map. The short exact sequence
	\begin{equation}\label{eq:gp}
		0 \to \Rp \to \Rq \to \Rq/\Rp \to 0
	\end{equation}
	thus shows that $g_\frakp \simeq \Sigma^{-1}(\Rq/\Rp)$. Moreover, \eqref{eq:gp} is a flat resolution of $\Rq/\Rp$. Since $\Rp \otimes_R k(\frakp) = k(\frakp)$ and $\Rq \otimes_R k(\frakp) = 0$ because an element of $\frakp \setminus \frakq$ is inverted in $\Rq$ but annihilates $k(\frakp)$, we obtain~\eqref{eq:gp-tensor-kp}. It follows that $k(\frakp) \in \Loco{g_\frakp}$ which is the inclusion $\supseteq$ in~\eqref{eq:gp-kp}.

	Now assume that $\frakp^2 \neq \frakp$. By~\cref{lem:p-idempotent-principal}, the maximal ideal of the valuation domain $\Rp$ is principal, say $\frakp\Rp = s\Rp$ with $s \in \frakp$, and its immediate generalization is $\frakq\Rp$. Applying~\cref{lem:for-one-for-all} to $\Rp$ gives $\Rp[1/s] = (\Rp)_{\frakq\Rp} = \Rq$. Hence $\Rq/\Rp = \bigcup_{n \ge 1} s^{-n}\Rp/\Rp$ and multiplication by $s^n$ identifies $s^{-n}\Rp/\Rp$ with $\Rp/s^n\Rp$, so that $\Rq/\Rp \cong \colim_n \Rp/s^n\Rp$ along the maps given by multiplication by $s$. Each~$\Rp/s^n\Rp$ is built from $\Rp/s\Rp = k(\frakp)$ by finitely many extensions, so it lies in $\Loco{k(\frakp)}$, and hence so does the colimit, by~\cref{lem:hocolim}. Therefore $g_\frakp \simeq \Sigma^{-1}(\Rq/\Rp) \in \Loco{k(\frakp)}$.
\end{proof}

\begin{Exa}
	Let $R$ be a rank 1 valuation domain that is not discrete. The maximal ideal~$\frakm$ is thus branched and $\frakm^2=\frakm$. The Balmer--Favi idempotents are $g_{(0)} = K$ and $g_\frakm \simeq \Sigma^{-1}(K/R)$ where $K=k(0)$ denotes the field of fractions. This follows from the proof of~\cref{prop:gp-kp} since the immediate generalization of~$\frakm$ is~$(0)$. Since $\frakm^2 = \frakm$, we have $\Tor_1^R(k(\frakm),k(\frakm)) = \frakm/\frakm^2 = 0$, so $k(\frakm) \otimes_R^{\bbL} k(\frakm) \simeq k(\frakm)$; that is, $R \to k(\frakm)$ is a smashing localization. Consequently, $\Loco{k(\frakm)}$ consists of those $X \in \Der(R)$ for which the map $X \to X \otimes_R^{\bbL} k(\frakm)$ induced by $R \to k(\frakm)$ is an isomorphism. If $g_\frakm$ were contained in $\Loco{k(\frakm)}$ then by~\eqref{eq:gp-tensor-kp} we would have $\Sigma^{-1}(K/R) \simeq g_\frakm \simeq g_\frakm \otimes_R^{\bbL} k(\frakm) \simeq k(\frakm)$ which is false since these complexes are concentrated in different degrees. Thus the second statement of~\cref{prop:gp-kp} fails for idempotent branched primes. In this example, the smashing localization $R \to k(\frakm)$ is not finite. Indeed, any perfect complex $P$ such that $P \otimes_R^{\bbL} k(\frakm) = 0$ has closed support not containing $\frakm$, so that $P = 0$. On the other hand, $\frakm$ is a nonzero acyclic object since $\frakm \otimes_R^{\bbL} k(\frakm) = \frakm/\frakm^2 = 0$.
\end{Exa}

\begin{Thm}\label{thm:tt-stratified}
	The derived category $\Der(R)$ of a valuation domain $R$ is stratified if and only if $R$ is strongly discrete.
\end{Thm}

\begin{proof}
	If $\Der(R)$ is stratified then $\Spec(R)$ is weakly noetherian by~\cite[Theorem~8.13]{Zou25} and $\Der(R)$ is homologically stratified by~\cite[Theorem~9.6]{BarthelHeardSandersZou26}. In particular,~$\Der(R)$ has the homological detection property. Thus $R$ is strongly discrete by~\cref{thm:h-detection-char} and~\cref{cor:strongly-discrete-noetherian}.

	Conversely, suppose $R$ is strongly discrete, that is, no nonzero prime ideal is idempotent. Then $\Spec(R)$ is noetherian by~\cref{cor:strongly-discrete-noetherian} and the equality~\eqref{eq:gp-kp} holds for all nonzero prime ideals by~\cref{prop:gp-kp}. It also holds for the zero ideal, since for the generic point $\eta$ we have $g_\eta = f_{\gen(\eta)^\complement} = R_\eta = K = k(\eta)$. Each of the localizing ideals $\Loco{g_\frakp}=\Loco{k(\frakp)}$ is minimal since $k(\frakp)$ is a field object. Furthermore, the local-to-global principle holds, either by~\cref{prop:LGP} or by~\cite[Theorem~3.22]{BarthelHeardSanders23a}. The result now follows from~\cite[Theorem~4.1]{BarthelHeardSanders23a}.
\end{proof}

The following corollary is originally due to Bazzoni--\Stovicek~\cite[Theorem~7.2]{BazzoniStovicek17}.

\begin{Cor}\label{cor:tc-char}
  Let $R$ be a valuation domain. The Telescope Conjecture holds for~$\Der(R)$ if and only if $R$ is strongly discrete.
\end{Cor}

\begin{proof}
	$(\Rightarrow)$: Suppose $\frakp \subset R$ is an idempotent prime ideal. Then $R/\frakp \otimes_R^\bbL \frakp = 0$ and
	\[
		\frakp \to R \to R/\frakp \to \Sigma \frakp
	\]
	is an idempotent triangle, i.e., $\Loco{\frakp}$ is a smashing ideal. If the telescope conjecture holds then this smashing ideal is compactly generated. Hence $\Supp^h(\frakp)$ is Thomason. However, $\frakp$ and $R/\frakp$ have disjoint h-support, so 
	\[
		\Supp^h(\frakp) = (\Supp^h(R/\frakp))^\complement = (\overbar{\{\frakp\}})^\complement.
	\]
	Hence $(\overbar{\{\frakp\}})^\complement$ is specialization closed. Since the prime ideals are totally ordered, $\frakp$ cannot have any proper generalizations, i.e., $\frakp = (0)$. We conclude that there are no nonzero idempotent prime ideals, that is, $R$ is strongly discrete.

	$(\Leftarrow)$: If $R$ is strongly discrete then $\Spec(R)$ is noetherian by \cref{cor:strongly-discrete-noetherian} and~$\Der(R)$ is stratified by \cref{thm:tt-stratified}. We can then invoke~\cite[Thm.~9.11]{BarthelHeardSanders23a}.
\end{proof}

We conclude this section with a clarification related to the proof of~\cref{prop:LGP}.

\begin{Rem}
	It was claimed in~\cite[Remark~7.8]{Zou25} that the spectrum $\Spc(\SHcp)$ of the category of finite $p$-local spectra satisfies the Hochster weakly scattered condition except for the empty Thomason subset. This claimed exception was an oversight. Indeed, in the notation of \cite[Example~6.5]{Zou25}, $\Spc(\SHcp)=\gen(\cat P_{p,\infty})$. Thus, for $Y=\emptyset$, taking $\cat P=\cat P_{p,\infty}$ and $U=\Spc(\SHcp)$, we obtain
	\[
		\cat P_{p,\infty} \in U \cap Y^\complement = \gen(\cat P_{p,\infty})
	\]
	which verifies the required condition. Thus, $\Spc(\SHcp)$ is Hochster weakly scattered. Consequently, the statement that $\SHp$ satisfies the local-to-global principle (\cite[Proposition~7.3]{Zou25}) also follows directly from~\cite[Theorem~7.6]{Zou25}.
\end{Rem}

\section{Homological stratification}

We now turn to characterizing when a valuation domain is homologically stratified in the sense of \cite{BarthelHeardSandersZou26}. This will take some preparation.

\begin{Rec}
	Let $A$ be a commutative ring. We say that $\Der(A)$ is \emph{homologically stratified} (or \emph{h-stratified}) if the map
	\[
		\{\text{localizing ideals of $\Der(A)$}\} \xrightarrow{\Supph} \{\text{subsets of $\Spec(A)$}\}
	\]
	is a bijection. This is equivalent to the \emph{homological local-to-global principle}:
	\begin{equation}\label{eq:h-LGP}
		A \in \Loco{k(\frakp) \mid \frakp \in \Spec(A)}.
	\end{equation}
	This equivalence is a consequence of the fact that $\Der(A)$ has enough homological residue fields; see \cite[Corollary~4.9 and Example~4.10]{BarthelHeardSandersZou26} for further details. Throughout this section we will use the characterization of h-stratification given by~\eqref{eq:h-LGP}. Note also that h-stratification implies the h-detection property. Indeed, if $\Supph(X) = \emptyset$, then the localizing subcategory $\SETT{Y}{Y \otimes_A^{\bbL} X = 0}$ contains every $k(\frakp)$, hence contains $A$ by~\eqref{eq:h-LGP}, so $X = 0$.
\end{Rec}

\begin{Cons}
	Let $R$ be a valuation domain. For $\frakp \subseteq \frakq$ in $\Spec(R)$, we may consider the valuation domain
	\[
		\Rpq \coloneqq R_\frakq/\frakp R_\frakq
	\]
	whose spectrum coincides with the closed interval $[\frakp,\frakq] \subseteq \Spec(R)$.
\end{Cons}

\begin{Def}
	We will say that $R$ is h-stratified on the interval $[\frakp,\frakq]$ if $\Der(\Rpq)$ is h-stratified.
\end{Def}

\begin{Rem}
	Note that $R_{[0,\frakq]} = R_{\frakq}$ and $R_{[\frakp,\frakm]} = R/\frakp$.
\end{Rem}

\begin{Lem}\label{lem:discrete-interval}
	Let $R$ be a valuation domain. The following two statements hold:
	\begin{enumerate}
		\item If $R$ is discrete then so is $\Rpq$ for any interval~$[\frakp,\frakq]$.
		\item If $\Rpq$ is discrete for each immediate inclusion $\frakp \subsetneq \frakq$ then $R$ is discrete.
	\end{enumerate}
\end{Lem}

\begin{proof}
	By \cref{rem:change-of-group}, the localization $R_\frakq$ has value group $\Gamma/H_\frakq$. Since $\frakp \subseteq \frakp R_\frakq \subseteq \frakp\Rp = \frakp$ by~\cref{lem:Rp-tensor-p}, we have $\frakp R_\frakq = \frakp$, so the convex subgroup of $\Gamma/H_\frakq$ associated to $\frakp R_\frakq$ is $H_\frakp/H_\frakq$. Applying \cref{rem:change-of-group} again, now to $R_\frakq \to R_\frakq/\frakp R_\frakq$, we see that $\Rpq$ has value group $H_\frakp/H_\frakq$. If $\frakp \subsetneq \frakq$ is an immediate inclusion, then $H_\frakq \subsetneq H_\frakp$ is an immediate inclusion of convex subgroups. Hence, by \cref{rem:discrete-iff-discrete} and \cref{def:discrete-group}, $R$ is discrete if and only if $\Rpq$ is discrete for each immediate generalization $\frakp \subsetneq \frakq$.  Moreover, for \emph{any} interval $[\frakp,\frakq]$, the convex subgroups of $H_\frakp/H_\frakq$ are the $H/H_\frakq$ with $H_\frakq \subseteq H \subseteq H_\frakp$. Thus, if $R$ is discrete then so is $\Rpq$ for any interval $[\frakp,\frakq]$.
\end{proof}

\begin{Not}
	In the following, we will write $k_A(\frakp)$ for the residue fields of a commutative ring $A$ whenever clarity is gained by specifying the ring. We will continue to simply write $k(\frakp)$ for the valuation domain~$R$.
\end{Not}

\begin{Lem}\label{lem:h-strat-restrict}
	Let $A \to B$ be a map of commutative rings and let $G\colon \Der(B) \to \Der(A)$ denote restriction of scalars. If $\Der(B)$ is h-stratified then 
	\[
		G(t) \in \Loco{G(k_B(\frakq)) \mid \frakq \in \Spec(B)}
	\]
	for every $t \in \Der(B)$.
\end{Lem}

\begin{proof}
	Since $\Der(B)$ is h-stratified, we have $t \in \Loco{k_B(\frakq) \mid \frakq \in \Spec(B)}$ by~\eqref{eq:h-LGP}.  The functor $G$ is exact and preserves coproducts, so the preimage under $G$ of $\Loco{G(k_B(\frakq)) \mid \frakq \in \Spec(B)}$ is a localizing ideal of $\Der(B)$. Since it contains each~$k_B(\frakq)$, it also contains~$t$. 
\end{proof}

\begin{Cor}\label{cor:h-strat-interval}
	Let $R$ be a valuation domain, let $\frakp \subseteq \frakq$ be prime ideals, set $B \coloneqq \Rpq$, and write $G\colon \Der(B) \to \Der(R)$ for restriction of scalars. For every $\mathfrak{s} \in [\frakp,\frakq]$ we have $G(k_B(\mathfrak{s}B)) \cong k(\mathfrak{s})$. Consequently, if $R$ is h-stratified on $[\frakp,\frakq]$ then \[ G(t) \in \Loco{k(\mathfrak{s}) \mid \mathfrak{s} \in [\frakp,\frakq]} \] for every $t \in \Der(B)$. In particular, this holds for $t = B$.
\end{Cor}

\begin{proof}
	Every prime ideal of $B = R_\frakq/\frakp R_\frakq$ is of the form $\mathfrak{s}B$ for a unique $\mathfrak{s} \in [\frakp,\frakq]$. Localizing $B$ at $\mathfrak{s}B$ gives $R_{\mathfrak{s}}/\frakp R_{\mathfrak{s}}$, compatibly with the maps from $R$, and its residue field is $R_{\mathfrak{s}}/\mathfrak{s}R_{\mathfrak{s}} = k(\mathfrak{s})$. This gives $G(k_B(\mathfrak{s}B)) \cong k(\mathfrak{s})$, and the rest follows from~\cref{lem:h-strat-restrict}.
\end{proof}

\begin{Lem}\label{lem:SES}
	For any prime ideal $\frakr \in \Spec(R)$, there is an exact sequence
	\[
		0 \to R \to R_\frakr \oplus R/\frakr \to k(\frakr) \to 0
	\]
	of $R$-modules where the last map is the difference of the two canonical maps.
\end{Lem}

\begin{proof}
	Injectivity of the first map follows from the injectivity of $R \to R_\frakr$ while surjectivity of the second map follows from the surjectivity of $R_\frakr \to k(\frakr)$. The composite vanishes since the second map is the difference of the canonical maps. Now suppose $a,b \in R$ and $s \notin \frakr$ are such that $(a/s,\overbar{b})$ maps to zero. Then $(a-bs)/s \in \frakr R_\frakr$, so $a - bs \in \frakr R_\frakr \cap R = \frakr$. Since $s \notin \frakr$ and the ideals of $R$ are totally ordered, $\frakr \subseteq (s)$, so $a - bs = cs$ for some $c \in R$. Moreover, $c \in \frakr$ since $cs \in \frakr$ and $s \notin \frakr$. Thus $a = (b+c)s$, and $b + c \in R$ maps to $(a/s, \overbar{b+c}) = (a/s, \overbar{b})$.
\end{proof}

\begin{Lem}\label{lem:glue}
	Let $\mathfrak{r} \in [\frakp,\frakq]$. If $R$ is h-stratified on $[\frakp,\mathfrak{r}]$ and on $[\mathfrak{r},\frakq]$ then $R$ is h-stratified on $[\frakp,\frakq]$.
\end{Lem}

\begin{proof}
	Replacing $R$ by $\Rpq$ and noting that $(\Rpq)_{\frakr} = R_{[\frakp,\frakr]}$ and $\Rpq/\frakr\Rpq = R_{[\frakr,\frakq]}$, the hypothesis becomes the statement that $R_\frakr = R_{[0,\frakr]}$ and $R/\frakr = R_{[\frakr,\frakm]}$ are h-stratified, and the goal is to show that $R$ is h-stratified. Applying \cref{cor:h-strat-interval} to the intervals $[0,\frakr]$ and $[\frakr,\frakm]$, with $t$ the ring itself in each case, we obtain
	\[
		R_\frakr \in \Loco{k(\mathfrak{s}) \mid \mathfrak{s} \in [0,\frakr]} \quad\text{and}\quad R/\frakr \in \Loco{k(\mathfrak{s}) \mid \mathfrak{s} \in [\frakr,\frakm]}
	\]
	in $\Der(R)$. The short exact sequence of \cref{lem:SES} then yields
	\[
		R \in \Loco{k(\mathfrak{s}) \mid \mathfrak{s} \in \Spec(R)}
	\]
	and we conclude that $\Der(R)$ is h-stratified.
\end{proof}

\begin{Lem}\label{lem:discrete-rank-one}
	If $R$ is a discrete rank one valuation domain then $\Der(R)$ is h-stratified.
\end{Lem}

\begin{proof}
	A discrete rank one valuation domain is noetherian and thus is h-stratified by \cite{Neeman92a} and \cite[Theorem 9.6]{BarthelHeardSandersZou26}.
\end{proof}

\begin{Lem}\label{lem:discrete-finite}
	If $R$ is discrete and has only finitely many prime ideals then the derived category~$\Der(R)$ is h-stratified.
\end{Lem}

\begin{proof}
	Note that if $R$ has only one prime ideal then it is a field and hence h-stratified. Otherwise, with \cref{lem:glue} in hand, it suffices to prove that $R$ is h-stratified on $[\frakp,\frakq]$ where $\frakp$ is an immediate generalization of $\frakq$. But in this case $\Rpq$ is a rank 1 discrete valuation domain and hence h-stratified by~\cref{lem:discrete-rank-one}.
\end{proof}

\begin{Lem}\label{lem:hocolim}
	Let $A$ be a commutative ring and let $M=\colim_{i \in I} M_i$ be a filtered colimit of $A$-modules. If $\cat L$ is a localizing subcategory of $\Der(A)$ with $M_i \in \cat L$ for each $i \in I$ then $M \in \cat L$.
\end{Lem}

\begin{proof}
    By \cite[Theorems~2.2 and 3.1]{AlonsoJeremiasSouto00}, localizing subcategories of the derived category of a Grothendieck abelian category are closed under filtered colimits of complexes.
\end{proof}

\begin{Lem}\label{lem:Rp-tensor-p}
	Let $R$ be a valuation domain and let $\frakp \in \Spec(R)$. We have an equality $\frakp\Rp =\frakp$ as $R$-submodules of the field of fractions $K$. Consequently, we have an isomorphism $\Rp \otimes_R^{\bbL} \frakp \simeq \frakp$ in $\Der(R)$.
\end{Lem}

\begin{proof}
	Since $\Rp$ is flat, $\Rp\otimes_R^{\bbL} \frakp \simeq \Rp\otimes_R \frakp = \frakp \Rp$. We have an evident inclusion $\frakp \subseteq \frakp\Rp$ in the field of fractions. Now let $a/s \in \frakp \Rp$ with $a \in \frakp$ and $s\not\in \frakp$. An inclusion $(s) \subseteq (a)$ would imply that $s \in \frakp$. Hence, $(a) \subseteq (s)$ since $R$ is a valuation domain. So $a = sb$ for some $b \in R$. Moreover since $a \in \frakp$ and $\frakp$ is prime we have $b \in \frakp$. Thus $a/s = sb/s = b$ is contained in $\frakp$. This establishes that $\frakp = \frakp \Rp$.
\end{proof}

\begin{Lem}\label{lem:below}
	Let $R$ be a valuation domain which is discrete. If $R_\frakp$ is h-stratified for every $\frakp \subsetneq \frakm$ then $R$ is h-stratified.
\end{Lem}

\begin{proof}
	If $\frakm$ is the only prime ideal then $R$ is a field and hence h-stratified. Suppose~$R$ is not a field. If $\frakm$ is branched then it has an immediate generalization $\frakp \subsetneq \frakm$. Then $R_{[\frakp,\frakm]}=R/\frakp$ is a rank one valuation domain which is discrete, hence is h-stratified by \cref{lem:discrete-rank-one}. On the other hand, $R_\frakp = R_{[0,\frakp]}$ is h-stratified by hypothesis. We conclude that $R$ is h-stratified by \cref{lem:glue}.

	It remains to treat the case where $\frakm$ is unbranched. Let 
    \[
		\cat L \coloneqq \Loco{k(\frakp) \mid \frakp \in \Spec(R)}.
    \]
    From the short exact sequence $0 \to \frakm \to R \to k(\frakm) \to 0$ we see that it suffices to prove that $\frakm \in \cat L$. Since $\frakm$ is unbranched, as an $R$-module $\frakm = \bigcup_{\frakp \subsetneq \frakm} \frakp$ is the union of the chain of nonmaximal prime ideals. By~\cref{lem:hocolim}, we are reduced to proving that $\frakp \in \cat L$ for each $\frakp \subsetneq \frakm$. Since $\Rp$ is h-stratified by hypothesis, \cref{cor:h-strat-interval} applied to the interval $[0,\frakp]$ and the object $\frakp\Rp \in \Der(\Rp)$ shows that $\frakp\Rp$, regarded as an $R$-module, lies in $\Loco{k(\mathfrak{s}) \mid \mathfrak{s} \in [0,\frakp]} \subseteq \cat L$. But $\frakp\Rp = \frakp$ by~\cref{lem:Rp-tensor-p}, which completes the proof.
\end{proof}

\begin{Lem}\label{lem:above}
	Let $R$ be a valuation domain which is discrete. If $R/\frakp$ is h-stratified for every $\frakp\neq 0$ then $R$ is h-stratified.
\end{Lem}

\begin{proof}
	Let $\cat L \coloneqq \Loco{k(\frakp) \mid \frakp \in \Spec(R)}$ and write $K=k(0)$ for the field of fractions. From the short exact sequence $0\to R\to K\to K/R \to 0$ we see that it suffices to prove that $K/R \in \cat L$.

	For any $0 \neq a \in R$, we can consider the $R$-submodule $a^{-1} R \subseteq K$ consisting of those fractions with denominator $a$. Note that if $a \mid b$ then $a^{-1}R \subseteq b^{-1}R$ and we have a filtered union $K = \colim_{0 \neq a} a^{-1} R$ ordered by divisibility. Since filtered colimits are exact, the short exact sequences $0 \to R \to a^{-1} R \to a^{-1}R/R \to 0$ yield 
	\[
		0 \to R \to K \to \colim_{0 \neq a} a^{-1}R/R \to 0.
	\]
	In other words, $K/R \simeq \colim_{a\neq 0} a^{-1}R/R$. Also note that the map $a^{-1} R \to R$ given by multiplication by $a$ induces an isomorphism $a^{-1}R/R \xrightarrow{\sim} R/aR$. In summary:
	\[
		K/R \simeq \colim_{a \neq 0} a^{-1}R/R \simeq \colim_{a \neq 0} R/aR.
	\]
	Since this is a filtered colimit, we are reduced (by \cref{lem:hocolim}) to showing that $R/aR \in \cat L$ for each nonzero nonunit $a \in R$. 

    Let $a \in R$ be a nonzero nonunit. By \cref{lem:radical-is-prime}, $\frakq \coloneqq \sqrt{aR}$ is a nonzero prime ideal. By assumption $R/\frakq$ is h-stratified, so \cref{cor:h-strat-interval}, applied to the interval $[\frakq,\frakm]$, shows that every $R/\frakq$-module lies in $\Loco{k(\mathfrak s) \mid \mathfrak s \in [\frakq,\frakm]} \subseteq \cat L$. We claim that it suffices to show that $\frakq^N \subseteq aR$ for some $N \ge 1$. Indeed, $\frakq^N$ then annihilates~$R/aR$, and hence the $\frakq$-adic filtration $R/aR \supseteq \frakq (R/aR) \supseteq \cdots \supseteq \frakq^N (R/aR) = 0$ is finite. Since the subquotients are all $R/\frakq$-modules, $R/aR$ has a finite filtration all of whose subquotients lie in $\cat L$. Hence $R/aR \in \cat L$ itself. 

    The prime $\frakq$ is branched by~\cref{lem:branch-equiv}, hence not idempotent because $R$ is discrete. By~\cref{lem:p-idempotent-principal} there is an $s \in \frakq$ with $\frakq \Rq = s \Rq$, and $\frakq \Rq = \frakq$ by~\cref{lem:Rp-tensor-p}. Hence $\frakq^2 = \frakq \cdot s\Rq = s\,\frakq\Rq = s\frakq$, and inductively $\frakq^{n+1} = s^n \frakq \subseteq s^nR$ for all $n \ge 0$. Since $s \in \sqrt{aR}$ we have $s^n \in aR$ for some $n$, and hence $\frakq^{n+1} \subseteq aR$. This completes the proof.
\end{proof}

These lemmas have all been building towards the next result. The final ingredient is the material on scattered spaces from~\cref{sec:scattered}.

\begin{Thm}\label{thm:one-direction}
	Let $R$ be a valuation domain that is discrete and whose totally ordered set of prime ideals is scattered. The derived category $\Der(R)$ is h-stratified.
\end{Thm}

\begin{proof}
	We prove, by transfinite induction on $\beta$, that $\Der(R)$ is h-stratified for every valuation domain $R$ which is discrete and whose spectrum is scattered of Cantor--Bendixson height $\beta$. So let $R$ be such a ring and let $X = \Spec(R)$ be the set of prime ideals. The hypothesis implies that~$X$ is a scattered space when equipped with the order topology (\cref{lem:totally-scattered}) or, equivalently, with its constructible topology (\cref{lem:constructible}). From either point of view, this space is quasi-compact. Hence by~\cref{prop:CB-successor}, $\CBheight(X) = \alpha + 1$ is a successor ordinal and $\delta^\alpha X$ is finite.
  
	We may write 
	\begin{equation}\label{eq:finite-delta}
		\delta^\alpha X = \{ \frakp_1 < \frakp_2 < \cdots < \frakp_n \}.
	\end{equation}
	If $\alpha = 0$ then $X=\delta^\alpha X$ is itself finite and hence the result is established by~\cref{lem:discrete-finite}. Thus, suppose $\alpha > 0$. 


	Before proceeding, we make the following general observation. For any inclusion of prime ideals $\frakp \subseteq \frakq$, the valuation domain $\Rpq$ is discrete (\cref{lem:discrete-interval}) and its spectrum, the interval $[\frakp,\frakq]$, is scattered. Moreover, since $[\frakp,\frakq]$ is convex in $X$, its own order topology agrees with its subspace topology. Thus, the Cantor--Bendixson height of $\Spec(\Rpq)$ is that of the subspace $[\frakp,\frakq] \subseteq X$. In particular, the inductive hypothesis applies to $\Rpq$ whenever this height is smaller than that of $X$.

	We now return to~\eqref{eq:finite-delta} with $\alpha > 0$. Note that for every $\frakp < \frakp_1$, we have $[0,\frakp]\cap \delta^\alpha X = \emptyset$. Hence the Cantor--Bendixson height of $[0,\frakp]$ is strictly smaller than that of $X$ by \cref{prop:drop-degree}. The inductive hypothesis implies that the valuation domain $R_{[0,\frakp]}$ is h-stratified. Since this is true for all $\frakp < \frakp_1$, \cref{lem:below} implies that $R_{[0,\frakp_1]}$ is h-stratified.

	Similarly, for every $\frakq > \frakp_n$, we have $[\frakq,\frakm] \cap \delta^\alpha X = \emptyset$. Thus $[\frakq,\frakm]$ has smaller Cantor--Bendixson height than $X$ and hence $R_{[\frakq,\frakm]}$ is h-stratified by the inductive hypothesis. Since this is true for all $\frakq>\frakp_n$, \cref{lem:above} implies that $R_{[\frakp_n,\frakm]}$ is h-stratified. By \cref{lem:glue} and ordinary induction, it remains to establish that~$R_{[\frakp_i,\frakp_{i+1}]}$ is h-stratified for each $i=1,\ldots,n-1$.
  
	If $[\frakp_i,\frakp_{i+1}] = \{\frakp_i,\frakp_{i+1}\}$, i.e., if $\frakp_i$ is an immediate generalization of $\frakp_{i+1}$, then we can apply~\cref{lem:discrete-rank-one}. Otherwise, we may choose a prime ideal $\frakp_i < \frakp' < \frakp_{i+1}$ and argue in a similar fashion as above: for every $\frakp_i < \frakq' \le \frakp'$, the interval $[\frakq',\frakp']$ does not intersect $\delta^\alpha X$, hence $R_{[\frakq',\frakp']}$ is h-stratified by the same inductive argument. Hence $R_{[\frakp_i,\frakp']}$ is h-stratified, again by invoking~\cref{lem:above}. Similarly, $R_{[\frakp',\frakr']}$ is h-stratified for all $\frakp' \le \frakr' < \frakp_{i+1}$. Consequently, $R_{[\frakp',\frakp_{i+1}]}$ is also h-stratified by~\cref{lem:below}. Again invoking~\cref{lem:glue}, we have that $R_{[\frakp_i,\frakp_{i+1}]}$ is h-stratified, which completes the proof.
\end{proof}

The final ingredient to our main theorem is the following recent result:

\begin{Thm}[Balchin--G\'{o}mez--Stevenson]\label{thm:BGS}%
  Let $A$ be a commutative ring. If $\Der(A)$ is homologically stratified then the constructible topology on $\Spec(A)$ is scattered.
\end{Thm}

\begin{proof}
	This is \cite[Theorem~5.3]{BalchinGomezStevenson26pp} bearing in mind that a topological space has Cantor--Bendixson rank (in their terminology) if and only if is scattered; see \cref{lem:scattered-CB}.
\end{proof}

\begin{Thm}\label{thm:h-stratification}
	Let $R$ be a valuation domain. The derived category $\Der(R)$ is homologically stratified if and only if $R$ is discrete and the ordered set $\Spec(R)$ contains no copy of $\bbQ$.
\end{Thm}

\begin{proof}
	Suppose that $\Der(R)$ is homologically stratified. Then $R$ is discrete by~\cref{thm:h-detection-char} and \cref{thm:BGS} implies that the constructible topology on $\Spec(R)$ is scattered. Hence the totally ordered set of prime ideals is scattered by~\cref{lem:constructible} and \cref{lem:totally-scattered}. In other words, it does not contain a suborder isomorphic to $\bbQ$. The converse is \cref{thm:one-direction}.
\end{proof}

\begin{Exa}
	Every strongly discrete valuation domain $R$ is homologically stratified. Indeed, by~\cite[Theorem~II.8.3]{FuchsSalce01}, $R$ is discrete and its prime ideals are well-ordered by reverse inclusion. In particular, $\Spec(R)$ is scattered and $\Der(R)$ is homologically stratified by~\cref{thm:h-stratification}. Alternatively, $\Der(R)$ is stratified by~\cref{thm:tt-stratified} and hence homologically stratified by~\cite[Theorem~9.6]{BarthelHeardSandersZou26}.
\end{Exa}

\begin{Cor}
	Let $R$ be a valuation domain whose value group $\Gamma=\prod_I^{\mathrm{lex}} \bbZ$ is a Hahn product of copies of the integers. The derived category $\Der(R)$ is h-stratified if and only if the totally ordered set $I$ contains no copy of $\bbQ$.
\end{Cor}

\begin{proof}
	The valuation domain $R$ is discrete by \cref{exa:hahn-discrete} and~\cref{rem:discrete-iff-discrete}. The ordered set of prime ideals is isomorphic to the set of convex subgroups of $\Gamma$ ordered by reverse inclusion (\cref{rem:convex}). This in turn is isomorphic to the set of final segments of $I$, ordered by reverse inclusion, by~\cref{prop:hahn-product-of-archimedean}. Initial segments are precisely the complements of final segments and complementation reverses inclusion. We thus conclude that the ordered set of prime ideals is isomorphic to the set of initial segments of $I$ ordered by inclusion. Thus $\Spec(R)$ is scattered if and only if the ordered set of initial segments of $I$ is scattered, which by~\cref{rem:initial-scattered}, is equivalent to $I$ itself being scattered. The result then follows from~\cref{thm:h-stratification}.
\end{proof}

\begin{Not}
	For any totally ordered set $I$, let $R_I$ denote any valuation domain whose value group is the Hahn product~$\prod_I^{\mathrm{lex}} \bbZ$. For example, $R_I$ could be the associated Hahn valuation domain of~\cref{cons:hahn-field}. Recall that $R_I$ is always discrete and that it is strongly discrete precisely when $I\op$ is well-ordered (\cref{rem:hahn-strongly-discrete}).
\end{Not}

\begin{Exa}
	The valuation domain~$R_{\bbQ}$ has the h-detection property (since it is discrete) but it is not h-stratified since the indexing set is $\bbQ$.
\end{Exa}

\begin{Exa}\label{exa:RNN}
	The valuation domain~$R_{\bbN}$ is h-stratified but not strongly discrete. This example is particularly interesting because it provides an example of an h-stratified category $\Der(R_{\bbN})$ which is not tt-stratified, as its spectrum is not weakly noetherian by~\cref{cor:strongly-discrete-noetherian}. It is also an example of an \mbox{h-stratified} category which does not satisfy the telescope conjecture by~\cref{cor:tc-char}. Finally, although it is not tt-stratified, it satisfies the local-to-global principle and hence the detection property. Note that the prime ideals of $R_\bbN$ form a chain
	\[
		0 = \frakp_0 \subsetneq \frakp_1 \subsetneq \frakp_2 \subsetneq \cdots \subsetneq \frakm
	\]
	with $\frakm = \bigcup_n \frakp_n$. Thus $\Spec(R_\bbN)\cong \Spc(\SHcp)$ is homeomorphic to the Balmer spectrum of the category of finite $p$-local spectra with the maximal ideal~$\frakm$ playing the role of the point at chromatic height~$\infty$. The categories $\Der(R_\bbN)$ and $\SHp$ thus share the same Balmer spectrum (of compact objects) but have different global properties. For example, the h-detection property does not hold for $\SHp$ due to the existence of nontrivial tensor-nilpotent objects, such as the Brown--Comenetz dual of the sphere. No such tensor-nilpotent objects can exist in $\Der(R_\bbN)$ since it \emph{does} have the h-detection property.
\end{Exa}

\section{Gabriel dimension}

The main theorem of this section (\cref{thm:gabriel-dimension}) is established in the classical literature, as will be explained in \cref{rem:classical}. However, understanding that proof requires juggling three different definitions of ``dimension'' and reading papers that are not easily accessible. We will provide a direct proof of the theorem.

\begin{Rem}
	The first ingredient is a characterization of strongly discrete valuation domains using ``transfinite powers'' of the maximal ideal:
\end{Rem}

\begin{Prop}\label{prop:jacobson-nilpotent}
	Let $R$ be a valuation domain with maximal ideal $J_0\coloneqq\frakm$. We define ideals $J_\alpha$ by transfinite recursion: $J_{\alpha + 1} \coloneqq \bigcap_{n \ge 1} J_\alpha^n$ for any ordinal $\alpha$ and $J_\lambda \coloneqq \bigcap_{\alpha < \lambda} J_\alpha$ for a limit ordinal $\lambda$. The following statements are equivalent:
	\begin{enumerate}
		\item $J_\alpha = 0$ for some ordinal $\alpha$;
		\item $R$ has no nonzero proper idempotent ideals.
	\end{enumerate}
\end{Prop}

\begin{proof}
	$(a)\Rightarrow (b)$: Suppose $I\subsetneq R$ is a proper idempotent ideal. We will prove that $I\subseteq J_\alpha$ for each ordinal $\alpha$. Firstly, we have $I \subseteq J_0 = \frakm$ since $I$ is proper. Now suppose $I \subseteq J_\alpha$. Then $I=I^2 \subseteq J_\alpha^2$ and, more generally, $I \subseteq J_\alpha^n$ for all $n \ge 1$. Hence $I \subseteq J_{\alpha+1} = \bigcap_{n \ge 1} J_\alpha^n$. For the limit step just observe that if $I \subseteq J_\alpha$ for all $\alpha < \lambda$ then $I \subseteq \bigcap_{\alpha < \lambda} J_\alpha = J_\lambda$. This completes the proof that $I \subseteq J_\alpha$ for all $\alpha$. Hence $I = 0$ by hypothesis.

	$(b)\Rightarrow (a)$: Since there is a proper class of ordinals and the ring~$R$ has only a set of ideals, the ideals $J_\alpha$ cannot all be distinct: There exists $\alpha < \beta$ such that $J_\alpha = J_\beta$. Moreover, since $\alpha + 1\le \beta$ we have $J_\alpha \supseteq J_{\alpha+1} \supseteq J_\beta = J_\alpha$, so that $J_\alpha = J_{\alpha+1}$. Consider this ideal~$I\coloneqq J_\alpha$. We have $I=J_{\alpha+1} = \bigcap_{n \ge 1} J_\alpha^n = \bigcap_{n \ge 1} I^n$ which implies that $I^2=I$. By construction $I \subseteq J_0 = \frakm$, so $I$ is proper. Hence $I=0$.
\end{proof}

\begin{Rec}
	Let $\cat A$ be a Grothendieck abelian category. The \emph{Gabriel filtration} of $\cat A$ is defined by transfinite recursion as follows. We set $\cat A_0$ to be the localizing subcategory generated by the simple objects\footnote{Some treatments state the definition using the finite length objects, but the simple objects and the finite length objects generate the same localizing subcategory.} of $\cat A$. For any ordinal $\alpha$ we define $\cat A_{\alpha + 1} \subseteq \cat A$ to be the localizing subcategory containing $\cat A_\alpha$ such that $\cat A_{\alpha + 1}/\cat A_\alpha = (\cat A/\cat A_{\alpha})_0$. Finally, for a limit ordinal $\lambda$ we set $\cat A_\lambda$ to be the localizing subcategory generated by $\bigcup_{\alpha < \lambda} \cat A_\alpha$. We say that $\cat A$ \emph{has Gabriel dimension} if $\cat A=\cat A_\alpha$ for some ordinal $\alpha$.
\end{Rec}

\begin{Rec}\label{rec:saturated}
	Let $\cat L \subseteq \cat A$ be a localizing subcategory and write $\cat R$ for the fully faithful right adjoint of the Gabriel localization $\cat Q\colon\cat A \to \cat A/\cat L$. We say  that an object $M \in \cat A$ is \emph{$\cat L$-torsion-free} if none of its nonzero subobjects are contained in~$\cat L$. We say that a subobject $N \subseteq M$ is \emph{$\cat L$-saturated} (or \emph{$\cat L$-pure}) if the quotient $M/N$ is $\cat L$-torsion-free.\footnote{This language comes from the hereditary torsion theory corresponding to the localization.} We remind the reader of two standard and readily verified facts:
	\begin{enumerate}
		\item If $N' \hookrightarrow Q(M)$ is any subobject then the pullback of $\cat R(N')\hookrightarrow \cat RQ(M)$ along the unit $M \to \cat RQ(M)$ yields an $\cat L$-saturated subobject $N \hookrightarrow M$ such that $Q(N)=N'$.
		\item The \emph{saturation} $\tilde{N}$ of a subobject $N \hookrightarrow M$ is defined to be the kernel of the composite $M \to M/N \to \cat RQ(M/N)$. It is an $\cat L$-saturated subobject of $M$ such that $\tilde{N}/N \in \cat L$. In particular, $Q(N)=Q(\tilde{N})$.
	\end{enumerate}
	These constructions provide an inclusion-preserving bijection between the subobjects of $Q(M)$ in $\cat A/\cat L$ and the $\cat L$-saturated subobjects of $M$ in $\cat A$.
\end{Rec}

\begin{Lem}\label{lem:cyclic-quotient}
	Let $R$ be a valuation domain, $\cat L \subseteq \Mod(R)$ a localizing subcategory, and $J \subseteq R$ an ideal. If $R/I \in \cat L$ for every $J \subsetneq I \subseteq R$ then $Q(R/J)\in\Mod(R)/\cat L$ is either zero or simple.
\end{Lem}

\begin{proof}
	A subobject of $Q(R/J)$ is the image of an $\cat L$-saturated subobject of $R/J$, hence is of the form $Q(I/J)$ for some $J \subseteq I \subseteq R$. If $J=I$ then the subobject is zero. Otherwise, we have $R/I \in \cat L$ by hypothesis. The short exact sequence
	\[
		0 \to I/J \to R/J \to R/I \to 0
	\]
	then shows that $Q(I/J) \simeq Q(R/J)$. We conclude that $Q(R/J)$ has at most two subobjects.
\end{proof}

\begin{Prop}\label{prop:jacobson-gabriel}
	Let $R$ be a valuation domain and let $\cat A\coloneqq\Mod(R)$ denote the category of $R$-modules. We have $R/J_\alpha \in \cat A_{\alpha + 1}$ for every ordinal $\alpha$.
\end{Prop}

\begin{proof}
	We prove this by transfinite induction. Note that $R/J_0 = R/\frakm$ is simple so $R/J_0 \in \cat A_0 \subseteq \cat A_1$. Thus the statement is true for $\alpha=0$.

	Now suppose $R/J_\alpha \in \cat A_{\alpha+1}$. Let $\beta = \alpha + 1$ and write $Q_\beta\colon \cat A \to \cat A/\cat A_\beta$ for the associated localization. Suppose we can prove that $R/I \in \cat A_\beta$ for every $J_\beta \subsetneq I \subseteq R$. Then~\cref{lem:cyclic-quotient} will imply that $Q_\beta(R/J_\beta) \in (\cat A/\cat A_\beta)_0$ so that $R/J_\beta \in \cat A_{\beta +1}$ as desired.

	With this in mind, consider any $J_\beta \subsetneq I \subseteq R$. By definition, $J_\beta = \bigcap_{n \ge 1} J_\alpha^n$. Let $x \in I$ be an element with $x \not\in J_\beta$ so that $x \not\in J_\alpha^n$ for some $n \ge 1$. Since ideals are totally ordered we must have $J_\alpha^n \subseteq xR$. Thus $J_\alpha^n \subseteq I$. The filtration $R \supseteq J_\alpha + I \supseteq J_\alpha^2 + I \supseteq \cdots\supseteq J_\alpha^{n-1} + I \supseteq I$ thus yields a finite filtration of $R/I$ whose subquotients are
	\[
		\frac{J_\alpha^{k-1} + I}{J_\alpha^{k}+I}
	\]
	for each $1 \le k \le n$, each of which is a quotient of $J_\alpha^{k-1}/J_\alpha^k$. In particular, each subquotient is annihilated by $J_\alpha$ and hence is an $R/J_\alpha$-module. This implies that each subquotient is contained in the localizing subcategory generated by $R/J_\alpha$ and hence is contained in $\cat A_{\beta}$ by hypothesis. We then conclude from the filtration that $R/I \in \cat A_{\beta}$. Since this is true for all $J_\beta \subsetneq I \subseteq R$ we conclude that $R/J_\beta \in \cat A_{\beta + 1}$.

	Now let $\lambda$ be a limit ordinal and suppose that $R/J_\alpha \in \cat A_{\alpha+1}$ for all $\alpha < \lambda$. By definition, $J_\lambda = \bigcap_{\alpha < \lambda} J_\alpha$. As before, consider any $J_\lambda \subsetneq I \subseteq R$. Choose an element~$x\in I$ which is not contained in $J_\lambda$. Thus $x \not\in J_\alpha$ for some $\alpha < \lambda$. This implies that $J_\alpha \subseteq I$ so we have a quotient $R/J_\alpha \twoheadrightarrow R/I$. Thus $R/I \in \cat A_{\alpha + 1}$ and note that $\cat A_{\alpha + 1} \subseteq \cat A_\lambda$ since $\alpha+1<\lambda$. Hence $R/I \in \cat A_\lambda$. Applying \cref{lem:cyclic-quotient} again, now with $\cat L = \cat A_\lambda$, we conclude that $Q_\lambda(R/J_\lambda)$ is zero or simple, that is, $R/J_\lambda \in \cat A_{\lambda+1}$.
\end{proof}

\begin{Prop}\label{prop:jacobson-not-gabriel}
	Let $R$ be a valuation domain and let $\cat A\coloneqq \Mod(R)$ denote the category of $R$-modules. If $J_\alpha \neq 0$ then $J_\alpha \not\in \cat A_{\alpha}$.
\end{Prop}

\begin{proof}
	Observe first that for any ideal $I$, the collection of $I$-power torsion modules
	\[
		\Mod_I(R) = \SETT{M \in \Mod(R)}{\text{for every } x \in M \text{ there exists } n \ge 1: I^n x =0}
	\]
	is a localizing subcategory of $\Mod(R)$. Indeed, closure under subobjects, quotients, and coproducts is immediate. For extension-closure, if $I^2 = I$, then $\Mod_I(R)$ consists of the modules annihilated by $I$ and is closed under extensions since $I^2 = I$. Otherwise, choose $a \in I \setminus I^2$. Since ideals are totally ordered, $I^2 \subseteq (a) \subseteq I$, and hence $I^{2n} \subseteq (a^n) \subseteq I^n$ for every $n \ge 1$. Thus $I$-power torsion coincides with $a$-power torsion, which is closed under extensions. We will prove:
	\begin{equation}\label{eq:A-Mod-J}
		\text{If $J_\alpha \neq 0$ then $\cat A_{\alpha} \subseteq \Mod_{J_\alpha}(R)$}.
	\end{equation}
	We will do this by transfinite induction. Note that $J_0 = \frakm$ and that the unique simple $R$-module is $k=R/\frakm$. We have 
	\[
		\cat A_0 = \Loc(R/\frakm) \subseteq \Mod_{\frakm}(R) = \Mod_{J_0}(R).
	\]
	For the inductive steps, keep in mind that if $\alpha < \beta$ then $J_\beta \subseteq J_\alpha$. Hence $\Mod_{J_\alpha}(R) \subseteq \Mod_{J_\beta}(R)$ and if $J_\beta \neq 0$ then $J_\alpha \neq 0$.

	With this in mind, suppose $J_{\alpha+1} \neq 0$ and that $\cat A_{\alpha} \subseteq \Mod_{J_\alpha}(R)$. Our main goal is to prove that every simple object in $\cat A/\cat A_\alpha$ is the image under ${Q:\cat A\to \cat A/\cat A_\alpha}$ of a $J_{\alpha+1}$-power torsion object. To this end, consider a simple object $S \in \cat A/\cat A_\alpha$. Since~$Q(R)$ is a generator, there is a nonzero morphism $Q(R) \to S$, which is necessarily surjective. Thus $S$ is the quotient of $Q(R)$ by a maximal subobject. By~\cref{rec:saturated} this maximal subobject corresponds to a maximal $\cat A_\alpha$-saturated subobject of $R$ in~$\cat A$. That is, we have $S = Q(R/I)$ for an ideal $I \subsetneq R$ which is $\cat A_\alpha$-saturated and maximal among all the proper $\cat A_\alpha$-saturated ideals of $R$.

	We claim that $J_{\alpha +1} \subseteq I$. Otherwise, there exists $x \in J_{\alpha+1}$ with $x \not\in I$. Hence $I \subsetneq xR$ since ideals are totally ordered. By the maximality of $I$ among $\cat A_{\alpha}$-saturated ideals, the $\cat A_{\alpha}$-saturation of $xR$ must be the whole of $R$, so that $R/xR \in \cat A_{\alpha}$. The inductive hypothesis says $\cat A_\alpha \subseteq \Mod_{J_\alpha}(R)$. Hence $\overbar{1} \in R/xR$ is annihilated by some power of $J_\alpha$, i.e., $J_\alpha^n \subseteq xR$ for some $n \ge 1$. On the other hand, $x \in J_{\alpha + 1} = \bigcap_{n \ge 1} J_\alpha^n$ so $xR\subseteq J_\alpha^n$. Hence $xR = J_\alpha^n$. Now observe that $x \in J_{\alpha+1} \subseteq J_{\alpha}^{2n} = x^2 R$ so that $xR = x^2 R$. Since $R$ is a domain, this implies $x$ is a unit, which is a contradiction since $x \in J_{\alpha+1} \subseteq J_0 = \frakm$. We conclude that $J_{\alpha+1} \subseteq I$. 

	This means that the ideal $J_{\alpha+1}$ annihilates the $R$-module $R/I$. In particular, $R/I \in \Mod_{J_{\alpha+1}}(R)$ and recall that $Q(R/I)=S$ is the simple object we started with. Now note that $J_{\alpha+1} \subseteq J_\alpha$ so $\cat A_{\alpha} \subseteq \Mod_{J_\alpha}(R) \subseteq \Mod_{J_{\alpha+1}}(R)$. Since every simple object in the quotient $\Mod(R)/\cat A_{\alpha}$ is contained in the image of $\Mod_{J_{\alpha+1}}(R)$, we conclude that $\cat A_{\alpha + 1} \subseteq \Mod_{J_{\alpha+1}}(R)$. This completes the successor step.

	Now let $\lambda$ be a limit ordinal and suppose $J_\lambda \neq 0$.  By the inductive hypothesis, we have $\cat A_\alpha \subseteq \Mod_{J_\alpha}(R)\subseteq \Mod_{J_\lambda}(R)$ for each $\alpha < \lambda$. Since $\Mod_{J_\lambda}(R)$ is a localizing subcategory, it follows that $\cat A_{\lambda} = \Loc(\bigcup_{\alpha < \lambda} \cat A_{\alpha}) \subseteq \Mod_{J_\lambda}(R)$.

	We have thus established~\eqref{eq:A-Mod-J} for all ordinals $\alpha$.

	Finally, note that if $x \in R$ is any nonzero element then $J_\alpha\neq 0$ implies $J_\alpha^n x \neq 0$ for all $n \ge 1$ because $R$ is a domain. Thus, $J_\alpha \neq 0$ implies that $\Mod_{J_\alpha}(R)$ does not contain \emph{any} nonzero ideal of $R$, hence neither does $\cat A_{\alpha}$ by~\eqref{eq:A-Mod-J}.
\end{proof}

\begin{Thm}\label{thm:gabriel-dimension}
	Let $R$ be a valuation domain. The category of $R$-modules $\Mod(R)$ has Gabriel dimension if and only if $R$ is strongly discrete.
\end{Thm}

\begin{proof}
	$(\Rightarrow)$: Suppose $\cat A=\Mod(R)$ has Gabriel dimension. Then $\cat A=\cat A_\alpha$ for some ordinal~$\alpha$. In particular, $J_\alpha \in \cat A_\alpha$. Thus~\cref{prop:jacobson-not-gabriel} implies that $J_\alpha = 0$. Hence~$R$ is strongly discrete by \cref{prop:jacobson-nilpotent}.

	$(\Leftarrow)$: Recall that we have $R/J_\alpha \in \cat A_{\alpha+1}$ for all ordinals $\alpha$ by \cref{prop:jacobson-gabriel}. If~$R$ is strongly discrete then it has no nonzero proper idempotent ideals by~\cref{lem:idempotent-is-prime} and hence $J_\alpha=0$ for some $\alpha$ by \cref{prop:jacobson-nilpotent}. Thus $R \in \cat A_{\alpha+1}$. Any nonzero $R$-module is a quotient of a coproduct of copies of $R$ and hence $\cat A=\cat A_{\alpha+1}$. This establishes that $\cat A$ has Gabriel dimension.
\end{proof}

\begin{Rem}\label{rem:classical}
As previously mentioned, \cref{thm:gabriel-dimension} is already established in the older literature on ring theory. Some clarifications on terminology will make this literature easier to navigate. Gabriel dimension was called Krull dimension in the original paper \cite{Gabriel62} and dubbed Gabriel dimension in \cite{GordonRobson74}. On the other hand, a Krull dimension for modules was introduced by \cite{RentschlerGabriel67} and this definition became ordinal-valued in~\cite{Krause70}. Lemonnier~\cite{Lemonnier72} then introduced the notion of ``deviation'' for partially ordered sets which, when applied to the lattice of submodules, is precisely this ordinal-valued Krull dimension. Be warned that the Krull dimension of a non-noetherian commutative ring, regarded as a module over itself, does not in general agree with its usual commutative algebra Krull dimension.

By \cite[Lemma~1.3(ii), Corollary~2.2 and Theorem~2.5]{GordonRobson74}, the category of modules $\Mod(R)$ has Gabriel dimension if and only if $R$ has Krull dimension in the sense of Rentschler--Gabriel--Krause \emph{and} every quotient of $R$ has finite uniform dimension (a.k.a.~Goldie dimension). Valuation domains are uniserial by definition, hence have uniform dimension $\le 1$. Hence, $\Mod(R)$ has Gabriel dimension if and only if $R$ has Krull dimension. Lemonnier~\cite[Proposition~4]{Lemonnier77} establishes that a valuation domain has Krull dimension if and only if it has no nonzero proper idempotent ideals, i.e., is strongly discrete. Although our proof using transfinite powers of the maximal ideal is not entirely trivial, we feel that it provides an interesting perspective on the theorem. The key transfinite construction is essentially a reordering of the definition given by Jacobson in \cite[Page~306]{Jacobson45}; see \cite{KrauseLenagan79}.
\end{Rem}

\bibliographystyle{alphasort}
\bibliography{bibliography}
\end{document}
